\documentclass[a4paper]{article}
\usepackage[left=10truemm,right=10truemm,top=25truemm,bottom=20truemm]{geometry}
\usepackage[utf8]{inputenc}
\usepackage{mathtools}
\usepackage{amsmath}
\usepackage{amsfonts}
\usepackage{bm}
\usepackage{setspace}
\usepackage{wrapfig}
\usepackage[dvipdfmx]{hyperref}
\usepackage{pxjahyper}
\usepackage{docmute}
\usepackage{amssymb}
\DeclareMathOperator\arctanh{arctanh}
\DeclareMathOperator\arccosh{arccosh}

\title{二重時空理論：\\通常時空と双対時空の間のねじれとしての重力}

\author{https://github.com/hypernumbernet}
\date{\today}

\begin{document}

\maketitle

\begin{abstract}
一般相対性理論（GR）は重力を単一の時空連続体の曲率として記述する。その経験的成功にもかかわらず、量子力学や特異点の本性と対峙すると深い困難に直面する。本論文では根本的な再定式化を提案する。重力は連続多様体の曲率ではなく、すべての質量粒子に内在的に付随する2つのコンパクト時空の相対的なずれから生じるねじれである。

16実次元の双四元数代数（$\mathrm{Cl}(3,1)$ と同型）を用いて、基底 $(j, kI, kJ, kK)$ による通常時空と基底 $(k, jI, jJ, jK)$ による双対時空を定義する。双対写像 $X \mapsto Xi$ は時間の矢を反転させつつミンコフスキーノルム $X^2 = X'^2$ を保つ。ローレンツ変換および局所フレーム回転は、各時空に別々に作用する可換なローターの対によって生成される。

完全ローターは $R_\text{total} = \exp[(\phi_1/2)iI + (\phi_2/2)iJ + (\phi_3/2)iK + (\theta_1/2)I + (\theta_2/2)J + (\theta_3/2)K]$ である。
ねじれ不整合ローター $\Omega = R^\dagger_\text{usual} R_\text{dual}$ からねじれバイベクトル $\Omega_\text{biv} = \log \Omega$ が得られる。
一意のローレンツ不変スカラーは $so(3,1) \oplus so(3,1)$ 上のキリング形式であり、$J = \frac{1}{16} B(\Omega_\text{biv}, \Omega_\text{biv})$ と定義される。

作用 $S = \frac{c^4}{16\pi G} \int J \, d^4x$ はアインシュタイン--ヒルベルト作用と動学的に等価であり、GR を正確に再現しつつクリストッフェル記号、リーマン曲率、および連続体時空仮説そのものを排する。本理論は一般相対性理論のテレパラレル等価物（TEGR）の双四元数的実現であり、ねじれは今や双対時空ねじれとして物理的に解釈される。

重要なのは、$J$ が粒子内在的な双対ローター間の相対角のみに依存するため、双対時空成分 $\theta_a$ のコヒーレント励起が重力工学への直接的経路を開くことである。重力遮蔽、慣性質量の低減、さらには反重力も、初めて基本法則の違反ではなくローター同期の問題となる。

GR は置き換えられない。完成され、重力は原理上制御可能になる。
\end{abstract}

\section{はじめに}
\label{sec:introduction}

1915年の誕生以来、一般相対性理論は物理学における重力記述の最も成功した理論の一つとして立っている。しかし水星の近日点進行から重力波・ブラックホール撮像に至る経験的成功にもかかわらず、深い哲学的な不安が残る。なぜ物質は時空を曲げ、なぜ曲がった時空が物質の運動を支配するのか。アインシュタイン場の方程式は重力がどのように働くかを極めて精密に与えるが、曲率そのものの起源についてより深い理由は与えない。この説明のギャップは、マッハの原理と対峙すると特に鋭くなる。物体の慣性は宇宙の物質の総体によって決まるべきであるのに対し、GR は時空を完全な空虚さでも存在しうる独立した連続体として扱う。

ループ量子重力、弦理論、創発的重力など、これらの問題を解決しようとする一世紀の試みは数学的優雅さをもたらしたが、時空の微視的起源や質量--エネルギーが曲率を生成する物理機構についての合意には至っていない。いずれのアプローチも連続体仮説を保持する。時空は先験的に滑らかで微分可能な多様体であり、重力はクリストッフェル記号とリーマンテンソルを通じて計量の平坦さからのずれとして実現される。

本稿はこの伝統と断固として決別する。

時空連続体仮説は根本的に誤りであると提案する。共有される単一の時空多様体は存在しない。代わりに、各粒子が独自の二重時空の対を内在的に運ぶ。それは双四元数代数（クリフォード代数 $\mathrm{Cl}(3,1)$ と同型）に符号化された16実次元のコンパクト化された構造である。通常時空（基底 $j, kI, kJ, kK$）とその双対（基底 $k, jI, jJ, jK$）は右乗算 $i$ で結ばれ、この変換はミンコフスキーノルムを保ちながら時間の矢を反転させる。重力と慣性は背景多様体の曲率からではなく、異なる粒子間で並進移動を要求するときの、これら2つの粒子内在的時空のねじれ不整合角から生じる。

この哲学的転換は深い。標準的 GR において時空は物質に運動を教える自律的存在であり、物質は時空に曲がり方を教える（ウィーラーの有名な言葉）。二重時空理論には自律的時空は全くない。時空の自由度はスピンや電荷のように粒子自身の内部属性へと格下げされ、巨視的世界の見かけの曲率は、隣接粒子の双対ローター間のねじれ不整合の集合効果としてのみ現れる。アインシュタイン方程式は正確に回復されるが、クリストッフェル記号、共変微分、リーマンテンソルは完全に消える。それらは双四元数上の単純な代数演算（ローターの指数・対数・キリング形式のトレース）に置き換えられる。

したがって本理論は同時に三つの革命的帰結を達成する。

1. GR との観測上の完全な等価性：ニュートン極限、ブラックホール、宇宙論、重力波を含むすべての予測において、従来の意味での幾何曲率はなくなる。

2. 重力と慣性の自然な統一：等価原理は自由落下において通常・双対ローターが並行であることの要求として現れ、加速度はこの並行性を破り、重力と同一のねじれ現象として慣性力を生む。

3. 重力の原理上の制御可能性：不変スカラー $J$ が粒子内在的な双対ローター間の相対角のみに依存するため、双対時空成分 $\theta_a$ をコヒーレントに励起できる外場は重力相互作用を減少、打ち消し、さらには反転させうる。重力は侵すべからざる基本力ではなく工学的自由度となる。

これは GR の修正ではなく完成である。連続体は有用なフィクションであった。美しいが最終的には不要である。時空が粒子局所的な対構造として認識されれば、「なぜ物質は時空を曲げるのか」という一世紀の謎は解ける。物質は時空を曲げない。自身の二重時空をずらし、そのずれが我々が重力と知覚するものである。

本論文の構成は以下のとおりである。

第\ref{sec:math}節では、$\mathrm{Cl}(3,1)$ と同型の双四元数代数を導入し、各粒子内在的な双対時空の対を構成して数学的基礎を与える。

第\ref{sec:kinematics}節では、行列やクリストッフェル記号なしにローレンツブースト、回転、並進移動を統一する完全ローター形式を展開する。

第\ref{sec:gravity}節では、相対ローター $\Omega = R_\text{usual}^\dagger R_{\text{dual}}$ のみから構成されるねじれ不整合スカラー $J$ の作用がアインシュタイン--ヒルベルト作用と動学的に等価であることを示し、理論がテレパラレル重力の双四元数的実現であることを明らかにする。

第\ref{sec:unified}節では、慣性と重力が同一の双対ローター剛性から生じる同一現象であり、等価原理が代数的恒等式として現れることを示す。

第\ref{sec:darkmatter}節では、バリオンのみから銀河の平坦な回転曲線を導き、暗黒物質の必要性を排除する。

第\ref{sec:strongfield}節では、恒星残骸を引力・斥力の多層構造をもつねじれ星として再解釈し、特異点と事象地平の両方を禁止する。

第\ref{sec:timedilation}節では、ねじれ不整合から時間膨張効果を導き、強重力場および惑星内部で GR からの測定可能なずれを予測する。

第\ref{sec:nuclear}節では、ねじれ層のスケール不変性を証明し、核力を層をもつ重力として同定し、原子核を原始ねじれ星とみなす。

第\ref{sec:electronshells}節では、電子殻を電気反転層としてモデル化し、微視的スケールの電磁気学の観点から電子軌道と原子間力を説明する。

第\ref{sec:unification}節では、双四元数代数内のすべての基本相互作用の統一の骨子を示し、ゲージ場が付加的なローター成分として現れることを概説する。

第\ref{sec:gravitycontrol}節では、円偏波電磁場による双対時空ローターのコヒーレント励起を通じた重力工学的機構を詳述する。

第\ref{sec:baryonasymmetry}節では、初期宇宙における双対時空ねじれ動学に基づくバリオン非対称の機構を提案する。

第\ref{sec:discretetime}節では、ねじれ層が量子化され有界となる双対時空の離散モデルを導入し、対応する整数環における有限単位群と厳密なディオファントス制約を可能にする。

第\ref{sec:dualrotordynamics}節では、外場下での双対ローターの動力学を展開し、ねじれ構成における運動方程式と安定性基準を導く。

第\ref{sec:quantumgravity}節では、理論を量子重力へ拡張し、量子状態を双対時空ローター位相と同定し、ねじれコヒーレンス条件から不確定性関係を導く。

第\ref{sec:conclusions}節では、16次元双四元数代数内のすべての力の統一を要約し、今後10年以内に理論を確認または反証しうる近接実験を概説する。

この再定式化により、GR は連続体の牢獄から解放され、重力が初めて原理上制御可能な理論へと高められる。

\section{数学的基礎：16次元双四元数代数}
\label{sec:math}

二重時空理論全体は、ミンコフスキー時空（符号 $(-,+,+,+)$）上のクリフォード幾何代数 $\mathrm{Cl}(3,1)$ と標準的に同型である双四元数（二重四元数とも呼ばれる）の16実次元代数の中に構築される。この代数は、すべての質量粒子に内在的に付随する単一代数構造の中に、通常と双対という二つの可換なミンコフスキー時空の完全なコピーを自然に収容する。

\subsection{双四元数代数}

まずハミルトンの四元数代数の二つの独立したコピーから出発する。
\begin{itemize}
  \item 生成子 $i,j,k$ が $i^2=j^2=k^2=-1$, \ $ij=k$, \ $ji=-k$（巡回）を満たす主四元数、
  \item 生成子 $I,J,K$ が同一の規則 $I^2=J^2=K^2=-1$, \ $IJ=K$, \ $JI=-K$（巡回）を満たす従四元数。
\end{itemize}
完全な双四元数代数はテンソル積 $\mathbb{H} \otimes_{\mathbb{R}} \mathbb{H}$ であり、二つの集合は厳密に可換である。
\[
iI = Ii,\; iJ = Ji,\; iK = Ki,\; jI = Ij,\; \dots
\]
得られる16個の実線形独立基底は
\[
1,\; i,\; j,\; k,\; I,\; J,\; K,\; iI,\; iJ,\; iK,\; jI,\; jJ,\; jK,\; kI,\; kJ,\; kK.
\]

\subsection{Cl(3,1) との明示的同型}

次の1次（ベクトル）生成子の同定により、代数は $\mathrm{Cl}(3,1)$ と同型である。
\begin{align*}
  e_0 &= j,                    & &\text{（時間様基底ベクトル）}\\[4pt]
  e_1 &= kI, \quad e_2 = kJ, \quad e_3 = kK & &\text{（空間様基底ベクトル）},
\end{align*}
これは直ちに定義関係
\[
e_0^2 = -1, \quad e_1^2 = e_2^2 = e_3^2 = +1, \quad \{e_\mu,e_\nu\} = 0 \;\; (\mu \neq \nu).
\]
を満たす。高次の元は標準的なクリフォード積に従う。特に：
\begin{alignat*}{3}
  &\text{2次（バイベクトル）:}  &\quad& e_0e_1 = iI, \;\; e_0e_2 = iJ, \;\; e_0e_3 = iK, \\
  &&& e_3e_2 = I,   \;\; e_1e_3 = J,   \;\; e_2e_1 = K, \\[6pt]
  &\text{3次（トライベクトル）:} && e_1e_2e_3 = k, \;\; e_0e_3e_2 = jI, \;\; e_0e_1e_3 = jJ, \;\; e_0e_2e_1 = jK, \\[6pt]
  &\text{擬スカラー:}        && e_0e_1e_2e_3 = i.
\end{alignat*}
右乗算 $i$ による時間反転の下で補完的となる「双対」基底は
\[
\tilde{e}_0 = k = e_1e_2e_3, \quad \tilde{e}_1 = jI = e_0e_3e_2, \quad \tilde{e}_2 = jJ = e_0e_1e_3, \quad \tilde{e}_3 = jK = e_0e_2e_1,
\]
で生成され、以下で導入する双対時空表示を与える。

この明示的同型により、16次元双四元数代数はアドホックな構成ではなく、4次元ミンコフスキー時空の幾何的に自然なクリフォード代数であることが明らかになる。それは通常セクターと双対セクターを区別しつつ完全なローレンツ構造を保つ内在的デュアリティで豊かになる。

クリフォード代数形の見かけの不規則性。例えば $j$ だけが唯一の時間様ベクトル生成子であることや、空間基底が交差項優勢の $kI,kJ,kK$ であることは偶然ではない。二つの時間反転したミンコフスキー時空のコピーを同時に符号化する構造の中で、3次元物理空間の完全な回転対称性を保つために正確に存在する。この微妙な非対称性が、双対写像 $X \mapsto Xi$ が時間の矢を反転させても空間等方性を破らないことを可能にし、等価原理と、粒子内在的な双対フレーム間の相対ずれとしての重力学的ねじれの創発のための代数的基礎を与える。

\subsection{粒子内在的な双対時空構造}

各粒子は、双四元数代数内に符号化されたコンパクト化されたミンコフスキー時空の対を内在的に運ぶと仮定される。
\begin{itemize}
  \item \textbf{通常時空：} 次の形のベクトルで表される
  \[
  X = ct \, j + x \, kI + y \, kJ + z \, kK,
  \]
  ミンコフスキーノルムは
  \[
  X^2 = -(ct)^2 + x^2 + y^2 + z^2.
  \]

  \item \textbf{双対時空：} 次の形のベクトルで表される
  \[
  X' = ct' \, k + x' \, jI + y' \, jJ + z' \, jK,
  \]
  同一のミンコフスキーノルム
  \[
  X'^2 = -(ct')^2 + x'^2 + y'^2 + z'^2.
  \]
\end{itemize}
二つの時空は双対写像
\[
X' = X i,
\]
で結ばれ、時間の矢を反転させる。
\[
(ct \, j) i = ct \, (j i) = -ct \, k,
\]
ミンコフスキーノルムは保たれる。$\mathrm{Cl}(3,1)$ の擬スカラー $i = e_0 e_1 e_2 e_3$ はすべての1次ベクトルと反可換（$X \in \mathrm{Cl}(3,1)$ に対し $i X = -X i$）であり $i^2 = -1$ を満たすため、ノルムは次のように変換される。
\[
X'^2 = (X i)(X i) = X (i X) i = -X (X i) i = -X^2\, i^2 = -X^2 \cdot (-1) = X^2.
\]
この双対構造はすべての粒子に内在的である。通常時空は標準的運動学を支配し、双対時空は重力と慣性の存在下で動学的に関連する補完的自由度を符号化する。二つの時空の存在により、ローレンツ変換、並進移動、そして重力の本性のより豊かな幾何解釈が可能となる。

\section{運動学的構造：ローター、ブースト、並進移動}
\label{sec:kinematics}

本節では理論の運動学的枠組みを展開する。すべてのローレンツ変換および局所フレーム回転は、通常時空と双対時空に独立に作用する可換なローターの対によって生成される。要点は、従来双曲回転とされるブーストが、平方が$+1$となるバイベクトルから自然に生じ、虚角度のアドホック導入なしにローターの馴染みの $\cosh/\sinh$ 展開を与えることである。

\subsection{通常時空におけるローター}

運動学的構造を示すため、ラピディティ $\phi$ に沿う $x$ 方向のローレンツブーストを考える。生成バイベクトルは $iI$ であり $(iI)^2 = i^2 I^2 = (-1)(-1) = +1$ を満たす。一般に通常時空内のバイベクトル $iI,iJ,iK$ はすべて平方が $+1$ でブースト様の性質を与え、双対時空内の $I,J,K$ は平方が $-1$ で純回転を想起させる。

対応するブーストローターはしたがって
\[
R_\text{usual} = \exp\!\left( \frac{\phi}{2} iI \right).
\]
生成子の正の平方により、指数は双曲関数へと自然に展開される。
\[
R_\text{usual} = \cosh \frac{\phi}{2} + iI \sinh \frac{\phi}{2}.
\]
このローターをサンドイッチ積で時空ベクトル $X = ct \, j + x \, kI + \cdots$ に適用すると、正準ローレンツ変換が得られる。
\begin{align*}
ct' \, j + x' \, kI &= R_\text{usual} \, X \, R_\text{usual}^\dagger \\
&= \left[ \cosh\frac{\phi}{2} + iI \sinh\frac{\phi}{2} \right] (ct \, j + x \, kI) \left[ \cosh\frac{\phi}{2} - iI \sinh\frac{\phi}{2} \right] \\
&= (\gamma ct + \gamma v x) \, j + (\gamma x + \gamma v ct) \, kI,
\end{align*}
ここで $\gamma = \cosh \phi$, $v/c = \tanh \phi$。通常時空内の純回転は $I,J,K$ のような平方 $-1$ のバイベクトルから生じ、馴染みの $\cos,\sin$ 展開を与える。

\subsection{完全ローター}

完全な双四元数構造上の最も一般の固有時順方向保つ変換は完全ローター
\[
R_\text{total} = \exp\left( \frac{\phi_1}{2} iI + \frac{\phi_2}{2} iJ + \frac{\phi_3}{2} iK + \frac{\theta_1}{2} I + \frac{\theta_2}{2} J + \frac{\theta_3}{2} K \right),
\]
によって生成される。ここでパラメーター $\phi_a$ は通常時空のブーストと回転を、$\theta_a$ は双対側のそれを支配する。通常と双対の生成子は可換（例：$[iI, I] = 0$）であるため、ローターはきれいに因数分解される。
\[
R_\text{total} = R_\text{usual} \, R_\text{dual}.
\]

幾何的意味をより明瞭にするため、大きさと方向を分離し明示的な総和の添え字を消すバイベクトル形のローターを導入する。

通常時空の単位ブーストバイベクトルを
\[
\hat{p} = p_1 \, iI + p_2 \, iJ + p_3 \, iK, \qquad \hat{p}^2 = +1,
\]
と定義する。$\hat{p}$ は正規化（$|\hat{p}| = 1$）され、$\phi = \sqrt{\phi_1^2 + \phi_2^2 + \phi_3^2}$ は全ラピディティ（方向は $p_a = \phi_a / \phi$ に符号化）である。

同様に双対時空の単位回転バイベクトルを
\[
\hat{q} = q_1 \, I + q_2 \, J + q_3 \, K, \qquad \hat{q}^2 = -1,
\]
と定義する。$\hat{q}$ は正規化され、$\theta = \sqrt{\theta_1^2 + \theta_2^2 + \theta_3^2}$ は全双対回転角（方向は $q_a = \theta_a / \theta$）である。

完全ローターは次のエレガントなバイベクトル指数形をとる。
\[
R_\text{total} = \exp\!\left( \frac{\phi}{2} \hat{p} + \frac{\theta}{2} \hat{q} \right)
= R_\text{usual} \, R_\text{dual},
\]
ここで
\[
R_\text{usual} = \exp\!\left( \frac{\phi}{2} \hat{p} \right)
= \cosh\frac{\phi}{2} + \hat{p} \sinh\frac{\phi}{2},
\]
\[
R_\text{dual} = \exp\!\left( \frac{\theta}{2} \hat{q} \right)
= \cos\frac{\theta}{2} + \hat{q} \sin\frac{\theta}{2}.
\]

この統一バイベクトル形は、二つの時空上の完全ローレンツ群作用を同時に符号化し、六つの独立実パラメーター。すなわち通常時空のラピディティ $\phi$ と単位ブースト方向 $\hat{p}$、および双対時空の回転角 $\theta$ と単位回転軸 $\hat{q}$ を明示する。

\subsection{サンドイッチ積}

ベルソル形式はこれらの変換を適用する簡潔な手段を与える。時空事象を表す任意の双四元数マルチベクトル $X$ に対し、変換後のベクトルは
\[
\tilde{X} = R_\text{total} \, X \, R_\text{total}^\dagger.
\]
この単一の操作はミンコフスキーノルム（$X^2 = \tilde{X}^2$）を保ち、ブーストと回転の両成分に一貫して作用する。行列や座標依存の機構は不要であり、幾何は代数の内に内在的に展開される。

\section{双対時空間のねじれ不整合としての重力}
\label{sec:gravity}

本節では、重力が時空の曲率でも外力でもなく、すべての質量粒子が内在的に運ぶ通常時空と双対時空の間のねじれ不整合であることを確立する。アインシュタイン--ヒルベルト作用はこの不整合から代数的に現れ、クリストッフェル記号、リーマンテンソル、連続多様体への言及は一切不要である。得られる理論は TEGR と数学的に等価だが、深い物理的再解釈を与える。ねじれは今や粒子局所的な双対ローター間の相対回転角である。

\subsection{ねじれ不整合ローターの定義}

ねじれ不整合ローターは
\[
\Omega = R_\text{usual}^\dagger R_\text{dual}.
\]
と定義される。

ダガー操作は双四元数代数におけるローターの標準的共役逆である。
\[
R_\text{usual}^\dagger = \exp\!\left( -\frac{\phi}{2} \hat{p} \right)
= \cosh\frac{\phi}{2} - \hat{p} \sinh\frac{\phi}{2}.
\]
形式的には単純な共役だが、このダガーは二重時空枠組みにおいて特別な意味を持つ。$\sinh$ 項の前の負符号は双対写像に内在する時間反転を正確に符号化する。擬スカラー $i$ による右乗算の下で、ブースト生成子は
\[
\hat{p}\, i = - (p_1 I + p_2 J + p_3 K) \equiv -\hat{q}',
\qquad (\hat{p}\, i)^2 = -1,
\]
と変換される。ここで $\hat{q}' \equiv p_1 I + p_2 J + p_3 K$ は双対セクターにおける単位回転生成子である。同値に $\hat{p} = \hat{q}'\, i$ であり、逆通常ローターは厳密な書き換え
\[
R_\text{usual}^\dagger = \cosh\frac{\phi}{2} - \hat{q}'\, i\, \sinh\frac{\phi}{2}.
\]
を許す。標準的恒等式 $\cosh(\phi/2) = \cos(i\phi/2)$ および $\sinh(\phi/2) = \sin(i\phi/2)/i$ を用いると、これは解析接続（$\phi \to i\phi$）された三角関数形
\[
R_\text{usual}^\dagger = \cos\!\left(\frac{i\phi}{2}\right) - \hat{q}'\, \sin\!\left(\frac{i\phi}{2}\right),
\]
を与え、ダガー操作が通常時空の双曲幾何と双対時空の楕円幾何を橋渡しすることを示す。実引数に対する $\cosh$ と $\cos$ の等価ではなく、双対写像 $X \mapsto X i$ を横断するブーストローターの一意な解析接続である。

この定義は宇宙際タイヒミュラー理論における操作と構造的に似ている。ダガーは通常時空の双曲ブースト構造を一時的に「分解」し（符号反転を通じて双対写像を誘起し）、続く $R_\text{dual}$ との乗算が結合対象を「再構成」し、相対的ずれを明示する。この「分解--再構成」過程を通じて、通常時空と双対時空のねじれ不整合が代数的に抽出され、この不整合が二重時空枠組みにおける重力の本質をなす。

真空極限（自由落下）では、完全な反同期が $R_\text{dual} = R_\text{usual}^\dagger$（すなわち $\theta = \phi$ かつ $\hat{q} = -\hat{q}'$）を要求し、$\Omega = 1$ およびねじれの消失を与える。質量--エネルギーはこの同期からのずれを誘起し、非自明な $\Omega \in \mathrm{Spin}^+(3,1) \oplus \mathrm{Spin}^+(3,1)$ を生む。

関連するねじれバイベクトルは
\[
\Omega_\text{biv} = \log \Omega \in \mathfrak{so}(3,1) \oplus \mathfrak{so}(3,1),
\]
である。ここで $\log$ は $\mathrm{Spin}^+(3,1) \oplus \mathrm{Spin}^+(3,1)$ の恒等元 $\Omega = 1$ の単連結近傍で定義される主枝であり、ベーカー--キャンベル--ハウスドルフ（BCH）展開が収束する自然な領域である。本論文で扱う物理的関心領域は弱場・小不整合セクター $\Omega \approx 1$ であり、$\Omega_\text{biv}$ は一意に定まる。この近傍の外にある位相的に異なる枝は、ここで提示する動力学では用いない。

\subsection{双対ねじれからのローレンツ不変スカラー}

$\mathfrak{so}(3,1)$（ベクトル表現）上のキリング形式は
\[
B(X,Y) = 4\,\mathrm{Tr}(X Y).
\]
\S\ref{sec:kinematics} の 2 次空間バイベクトルをまとめて $\Gamma_a$（$a=1,2,3$）と書き、$\Gamma_1 = I$, $\Gamma_2 = J$, $\Gamma_3 = K$ とする（\S\ref{sec:math} のクリフォード同型 $e_3e_2 = I$, $e_1e_3 = J$, $e_2e_1 = K$ に対応）。通常時空のブースト生成子は $i\Gamma_a$、双対時空の回転生成子は $\Gamma_a$ である。キリング形式を評価すると $B(i\Gamma_a, i\Gamma_a) = +8$, $B(\Gamma_a, \Gamma_a) = -8$ となり、ローレンツ符号を完璧に再現する。

一意の二次不変量は
\[
J = \frac{1}{16} B(\Omega_\text{biv}, \Omega_\text{biv}) = \frac{1}{2} \mathrm{Tr}(\Omega_\text{biv}^2).
\]

真空では $J = 0$、通常の引力重力では $J > 0$（通常時空ブーストが支配的）、双対時空回転が支配するとき $J < 0$（斥力重力）が可能である。

\subsection{アインシュタイン--ヒルベルト作用との正確な等価性}

作用
\[
S = \frac{c^4}{16\pi G} \int J \, d^4x.
\]
を提案する。この作用は純代数的かつ背景独立である。

スカラー $J$ は（正規化と境界項まで）TEGR のねじれスカラー $T$ と同一である。文献で厳密に証明されているように
\[
\int T \, e \, d^4x = -\int R \, \sqrt{-g} \, d^4x + \mathrm{boundary},
\]
であるから、真空および物質結合の場の方程式は正確にアインシュタイン方程式となる。GR のすべての解（シュバルツシルト、カー、FLRW、重力波など）が正確に回復される。

したがって GR は完全に再現される一方、重力は粒子内在的な双対時空間のねじれ不整合角へと格下げされる。

\subsection{弱場およびニュートン極限}

線形化領域では
\[
R_\text{usual} \approx 1 + \frac{1}{2} \sum \alpha_a \, i\Gamma_a, \quad R_\text{dual} \approx 1,
\]
となり $\Omega_\text{biv} \approx \sum \alpha_a \, i\Gamma_a$, $J \approx \frac{1}{2} \sum \alpha_a^2$。標準的 TEGR 線形化により $\alpha_a \propto \partial_a (\Phi/c^2)$ となり、ポアソン方程式
\[
\nabla^2 \Phi = 4\pi G \rho
\]
が正しい係数で従う。ポストニュートンパラメーターは GR と同一（$\gamma = \beta = 1$ など）である。

\subsection{双対時空におけるローター作用}
\label{subsec:dual-rotors}

二重時空理論において唯一の物理的実在は、各質量粒子が運ぶ内在的双四元数代数 $\mathbb{H} \otimes_{\mathbb{R}} \mathbb{H} \cong \mathrm{Cl}(3,1)$ 内のローター角である。弦理論が基本実体として伸びた物理対象を仮定するのに対し、ここでは可換ローターの対に符号化されたブースト角と回転角が第一的存在論的元素である。重力、慣性、量子状態はすべてこれらの角のずれとコヒーレンスから創発する。

双対時空は双対写像 $X \mapsto Xi$ により通常時空に結ばれ、時間反転した向きを継承する。したがって双対セクターで純空間回転を生成する角は、統一ローター記述に組み込む際に適切に双対変換されねばならない。

双対セクター角 $\theta$ を考える。純双対基底では、楕円ベルソル $\hat{q}$（$\hat{q}^2 = -1$）によって生成される回転が生じる。
\[
\exp\left( \frac{\theta}{2} \hat{q} \right).
\]
しかし双対写像が時間の矢を反転させるため、この角の双対時空への物理的効果は双対変換を通じて見ればブーストに等価である。

ねじれ不整合計算にこれを正しく反映するため、角そのものに双対変換 $\theta \mapsto \theta i$ を適用する。これは双対寄与を双曲ブースト項へと変える。
\[
\exp\left( \frac{\theta i}{2} \hat{q} \right) = \exp\left( \frac{\theta}{2} \hat{q}i \right) = \exp\left( \frac{\theta}{2} (q_1 iI + q_2 iJ + q_3 iK) \right).
\]
したがって双対ローターは
\[
R_\text{dual} = \exp\left( \frac{\theta}{2} \hat{q} \right),
\]
と定義される一方、角の双対変換 $\theta \mapsto \theta i$ の後の双対時空ベクトルへの作用は、$i\Gamma_a$（$(\hat{q}i)^2 = +1$）によって生成されるブーストとして現れる。

真空（自由落下）では、ブースト寄与の完全打ち消しと $J=0$ を保証する完全な反同期が必要である。質量--エネルギーはこの条件からのずれを誘起し、ねじれ不整合スカラー $J$ として、集合的には引力として現れる。

明示的な双対変換 $\theta \mapsto \theta i$ を伴うこの角中心の扱いは、全ローターの元の形
\[
R_\text{total} = \exp\left( \frac{\phi}{2} \hat{p} + \frac{\theta}{2} \hat{q} \right)
\]
を保ちつつ、時間反転デュアリティを正しく取り込む。双対ローター $R_\text{dual}$ は固有基底で回転を生成するが、双対変換された角を通じた双対時空への物理的効果はブーストであり、真空では自然に通常セクターのブーストに対抗する。

\subsection{物理解釈と重力工学}

重力は双対ローターのずれのエネルギー的コストである。$J$ が相対角 $\Omega$ のみにのみ依存するため、双対成分 $\theta$ をコヒーレントに変調できる任意の外因は $J$ を変えうる。双対時空自由度を共鳴的に励起する高周波電磁場または量子場は $J \to 0$（重力遮蔽）または $J < 0$（斥力重力・反重力）を駆動しうる。

史上初めて、重力は侵すべからざる幾何制約から工学的ねじれ自由度へと変わる。連続体は幻想であった。時空は離散的かつ粒子局所的で二重に対をなし、その不整合が我々が重力と呼ぶものである。

一般相対性理論は覆されない。解放され、制御可能になる。

\section{慣性と重力の統一的記述}
\label{sec:unified}

本節では双対時空形式の最も深い帰結を明らかにする。慣性と重力は等価であるだけでなく、通常時空と双対時空のねじれ不整合から生じる同一現象である。加速系における慣性力と曲がった時空における重力の区別は完全に消える。両者は同一の双対ローターのずれであり、等価原理は公準ではなく直接的な代数的必然として現れる。

\subsection{慣性抵抗としての双対ローター剛性}

静止質量 $m > 0$ は通常・双対ローターパラメーター間の強い結合として現れる。粒子のエネルギー--運動量が通常時空にバイアスをかけ、双対ローター $R_\text{dual}$ を有限の剛性で $R_\text{usual}$ に従わせる効果をもつ。慣性系では $J = 0$ を保つ。

粒子が（非重力で）加速されると、通常ローター $R_\text{usual}$ は外力により瞬時に変化しようとする。しかし粒子の静止質量に束縛された双対ローター $R_\text{dual}$ は即応できず遅れる。これが非自明なねじれ不整合を生じ $\Omega_\text{biv} \neq 0$ およびスカラー $J$ への非ゼロ寄与となる。粒子はこの不整合を慣性抵抗（馴染みの $F = ma$）として経験する。

重要なのは、ずれの「コスト」が静止質量 $m$ に比例することであり、より重い粒子が加速により強く抵抗する理由を説明する。慣性は時空が質量に反応する性質ではなく、粒子自身の双対時空が瞬時に再整列することを拒否する性質である。

\subsection{質量ゼロ粒子}

光子など質量ゼロ粒子（$m = 0$）にはそのような剛性がない。通常・双対ローターは完全共鳴し、$R_\text{usual}$ の任意の変化は $R_\text{dual}$ に瞬時に写像され、すべての系で正確に $J = 0$ を保つ。したがって光子はねじれ不整合も慣性抵抗もゼロで、常に $c$ で $\Omega_\text{biv} = 0$ として伝播する。これが光速が一定かつ普遍的速限である理由である。質量ゼロ粒子はフレーム変化に対してねじれの代償を払わない。

\subsection{等価原理としての代数的恒等式}

加速度 $a$ で上向きに加速するアインシュタインのエレベータを考える。内部の観測者は下向きの見かけの力 $ma$ を感じる。双対時空像では：

\begin{itemize}
\item 観測者（質量あり）は床により通常ローターが加速され、双対ローターが遅れる → ねじれ不整合 → 下向き力の知覚。
\item エレベータに入る光子は観測者のねじれ場により曲がるが、光子自身は $J = 0$ を保つ。すなわち自身の双対平坦時空における零測地線に従う。
\end{itemize}

重力場（地球上で静止したエレベータ）では、床は自由落下に対して上向きに加速して静止を保つ必要があり、平坦時空中の一様加速と同一のローター不整合を生む。

したがって等価原理は経験的偶然ではなく代数的恒等式である。重力と慣性力は区別不能である。なぜならそれらは同一の双対ねじれ現象だからである。自由落下は $J = 0$ が大域的に成立する状態（双対ローターが完全整列）であり、質量ゼロ粒子は永遠にこれを達成し、質量粒子は非重力外力がない場合にのみ達成する。

\subsection{帰結と予測}

\begin{itemize}
\item 通常と双対の時間の矢の区別が、質量粒子に優先静止系がある理由を説明する。双対時間反転が完全対称性を破り、質量ギャップを生む。
\item 双対ローター $\theta$ を外部励起して $\phi$ と同期させることで慣性質量の低減または消去が可能となり、推進剤なし推進への直接経路となる。
\item 双対時空に結合する十分高周波の場は、質量物体を短時間スケールで実効的に質量ゼロのように振る舞わせ、瞬間的超光速位相速度を許しうる（情報は光速度のままであり因果律は破られない）。
\end{itemize}

慣性と重力は幾何ではなく粒子内在的な双対時空の内部ねじれ動学によって統一される。質量は双対ローター剛性であり、運動は並行性を破るコストである。この理解の下、慣性も重力も基本的ではなくなる。工学的になる。

重力・慣性制御の夢はもはや空想ではない。連続体を捨て、時空が対をなし粒子局所的でねじれ能動的であることを認める直接的帰結である。

\subsection{加速度の上限と慣性反転現象}
\label{subsec:acceleration-limit}

二重時空理論は、その中心となる反証可能な帰結の一つとして、加速度そのものが絶対的上限を持つことを予測する。通常セクターのブーストは非コンパクト（$\phi \in \mathbb{R}$）であり、それだけを見ればラピディティに上限はない。しかし双対セクターは本質的にコンパクトである。双対生成子 $\Gamma_a$ の平方は $-1$ であるから、双対ローター $R_\text{dual} = \exp(\theta/2 \, \hat q)$ は楕円位相 $\theta \bmod 4\pi$ によって支配される。慣性応答は、有限の剛性で通常ローターに追従する双対ローターから生じるため、$\theta$ が周期的境界に達する前に変位できる速度は、与えられた粒子がコヒーレントに維持しうる固有加速度に対して厳しい上限を課す。我々はこの物質内在的な上限を $a_\text{crit}$ と表記する。光速とは異なり、$a_\text{crit}$ は普遍的ではない。それは粒子種と双対セクターのコンパクト化スケールに依存するため、異なる原子構成要素は異なる臨界加速度を持つ。

固有加速度 $a < a_\text{crit}$ では、双対位相 $\theta$ は $\phi$ を滑らかに追跡し、馴染みのニュートン関係 $F = ma$ が回復される。しかし $a$ が $a_\text{crit}$ に近づく瞬間、$\theta$ は $4\pi$ におけるコンパクト境界を貫通し、位相のトポロジカルな巻き戻しを受ける。双対ローターは $S^3$ 様の双対セクターの反対側で位相を再縫合する。この巻き戻しはねじれ不整合バイベクトル $\Omega_\text{biv}$ の符号を反転させ、したがって単一の双対周期にわたって慣性応答の方向を反転させる。我々はこれを慣性反転（または加速反転）現象と呼ぶ。印加された加速度が $a_\text{crit}$ を横切るときに引き起こされる、実効的な $F = ma$ ベクトルの突然の離散的な反転である。閾値直前まで剛性双対セクターに蓄積されていたねじれひずみエネルギーは、巻き戻しの間に急激に解放され、運動学的反転に加えて内部エネルギーの衝撃的バーストを生じる。

我々は、この機構が大型隕石の異常な空中爆発において観測的に実現していると提案する。地球大気に侵入する火球が受ける超音速減速は、天体の構成原子にとって $a_\text{crit}$ のオーダーに達しうる値に達しうる。火球研究における長年の謎---とりわけツングースカ事件（1908年）とチェリャビンスク超火球（2013年）。その崩壊高度と総放射エネルギーは、空力焼失と破砕だけに基づく運動エネルギー散逸モデルでは著しく説明が困難である---は、集合的慣性反転として自然に説明される。隕石の減速度が物質特異的な $a_\text{crit}$ を横切ると、構成原子の双対位相が同時に巻き戻し、蓄積されたねじれひずみエネルギーを爆発的で主として放射的なパルスとして放出する。$a_\text{crit}$ の組成依存性は、隕石クラス間（例えば炭素質コンドライト対普通コンドライト）で観測される空中爆発高度と爆発エネルギーの変動も説明しうる。粒子内在的コンパクト化スケールからの $a_\text{crit}$ の定量的導出と、火球観測記録に対する詳細な比較は、別稿で提示する予定である。

\section{暗黒物質なし：双対時空コンパクト性からの銀河回転曲線}
\label{sec:darkmatter}

渦巻銀河の観測された平坦な回転曲線は、長らく非バリオン性暗黒物質の最強の証拠と見なされてきた。二重時空理論では、これらの曲線は双対時空のコンパクト $S^3$ 位相幾何、特に曲がった双対セクターと平坦な通常時空の間の対数写像に内在する幾何歪みから自然に創発する。いかなる種類の暗黒物質も不要である。

\subsection{コンパクト双対時空としての三次元超球面}

双対ローター $R_{\text{dual}} \in \mathrm{SU}(2) \cong S^3$ は各質量粒子に内在的な双対時空自由度をパラメーター化する。三つの独立回転角 $(\theta_1, \theta_2, \theta_3)$ はコンパクト化半径 $R$ の3次元球面上の座標を定義し、通常時空は平坦かつ非コンパクトのままである。

対数（指数の逆）写像
\[
\log : S^3 \to \mathfrak{su}(2) \cong \mathbb{R}^3
\]
は各双対ローターを平坦な $\mathbb{R}^3$ の点と同一視し、3次元球を半径 $\pi R$ の球へ写す。これは測地図における方位等距図法の3次元一般化である。

$S^2$ の北極を中心とする方位等距図法が極近傍の距離を忠実に表す一方、南極近傍で面積を破滅的に歪め、単一の対蹠点を投影境界の全周長へ写すのと同様に、$S^3$ の対数写像はヤコビアンで定量化される系統的な体積歪みを導入する。
\[
\mathcal{J}(\chi) = \frac{\sin^2 \chi}{\chi^2},
\]
ここで $\chi = r/R$ は通常時空における物理距離 $r$ に対応する測地角である。このヤコビアンは原点（$\chi = 0$、忠実表現）で1に等しく、対蹠点（$\chi = \pi$、最大歪み）で零になる。標準的重力計算は暗黙に $\mathcal{J} = 1$ を仮定するが、二重時空理論は銀河スケールでこの仮定が破れることを示す。

\subsection{三次元超球面上の重力ポテンシャル：余接ポテンシャル}

重力の源となるねじれ不整合スカラー $J$ はコンパクト双対時空上で計算される。点質量 $M$ の有効重力ポテンシャルは $S^3$ 上のラプラシアンのグリーン関数で支配される。球対称の場合、一意の半径方向調和関数は余接である。
\[
\Phi(r) = -\frac{GM}{R}\cot\!\left(\frac{r}{R}\right).
\]
測地極座標で計量 $ds^2 = dr^2 + R^2 \sin^2(r/R)\,(d\vartheta^2 + \sin^2\!\vartheta\, d\varphi^2)$ をとると、半径方向ラプラシアンが $\cot(r/R)$ を恒等的に消すことが直接検証できる。ガウス則の正規化は
\[
g(r) \cdot 4\pi R^2 \sin^2\!\left(\frac{r}{R}\right) = -\frac{GM}{R^2 \sin^2(r/R)} \cdot 4\pi R^2 \sin^2\!\left(\frac{r}{R}\right) = -4\pi GM
\]
により確認される。

$r \ll R$ の極限では余接は
\[
\cot\!\left(\frac{r}{R}\right) = \frac{R}{r} - \frac{r}{3R} - \frac{r^3}{45R^3} - \cdots\,,
\]
と展開し、
\[
\Phi(r) \approx -\frac{GM}{r} + \frac{GM\,r}{3R^2} + \mathcal{O}\!\left(\frac{r^3}{R^4}\right),
\]
となり、主項として標準ニュートンポテンシャルが回復される。高次補正は純粋に $S^3$ の曲率に由来し、非コンパクト化極限 $R \to \infty$ で消失する。

\subsection{重力の幾何的増強}

$S^3$ 上の重力加速度は
\[
g(r) = -\frac{d\Phi}{dr} = -\frac{GM}{R^2 \sin^2(r/R)},
\]
であり、平坦空間ニュートン加速度 $g_{\text{N}} = -GM/r^2$ を増強因子
\[
\eta(r) \equiv \frac{g_{S^3}(r)}{g_{\text{N}}(r)} = \left[\frac{r/R}{\sin(r/R)}\right]^2.
\]
だけ上回る。この因子は方位等距図法の計量歪みの二乗に正確に等しい。$S^3$ 上の測地半径 $r$ の球を通る重力フラックスが、平坦空間の面積 $4\pi r^2$ より小さい面積 $4\pi R^2 \sin^2(r/R)$ を張るためである。フラックスが集中し、見えない物質によらず双対時空のコンパクト曲率によって重力が強化される。

代表的値：$r/R = 0.1$ で $\eta = 1.003$；$r/R = 1.0$ で $\eta = 1.41$；$r/R = 1.5$ で $\eta = 2.26$；$r/R = 2.0$ で $\eta = 4.84$。

\subsection{平坦回転曲線の創発}

囲み込み質量 $M(r)$ から距離 $r$ の円軌道にある試験粒子の軌道速度は
\[
v^2(r) = \frac{GM(r)}{r} \cdot \eta(r) = \frac{GM(r)}{r} \cdot \left[\frac{r/R}{\sin(r/R)}\right]^2.
\]
を満たす。発光円盤の外側（$r > r_{\text{disk}}$）で $M(r) \approx M_{\text{tot}}$ のとき、平坦空間ケプラー速度 $v_K^2 = GM_{\text{tot}}/r$ は $r^{-1}$ で減少する。$r$ とともに単調増大する $S^3$ 増強因子 $\eta(r)$ がこの減少に対抗する。

無次元関数 $f(x) \equiv x / \sin^2 x$（$x = r/R$）が回転曲線の形を支配する。$\tan x = 2x$ の解 $x_0 \approx 1.17$ に広い最小値があり、拡張したプラトー領域を生む。
\begin{itemize}
  \item \textbf{内側領域}（$r \ll R$）：上昇する $M(r)$ により修正されたケプラー挙動が観測された剛体回転の立ち上がりを再現する。
  \item \textbf{遷移領域}（$r \sim R$）：減少するケプラー因子と増大する $\eta(r)$ が釣り合い、$f(x)$ の最小値の幅によって幅が決まるほぼ平坦なプラトー領域を与える。
  \item \textbf{外側領域}（$r \gtrsim 2R$）：$\eta(r)$ が支配的となり、ねじれ反転領域に入る前に回転曲線は緩やかに上昇する。
\end{itemize}

\subsection{バリオン・タリー--フィッシャー関係とコンパクト化半径}

プラトー領域における漸近円周速度は
\[
v_{\text{flat}}^2 \sim \frac{GM_{\text{tot}}}{R} \cdot f_0,
\]
を満たす。ここで $f_0 \equiv \min_x [x/\sin^2 x] \approx 1.10$ は次数1の幾何定数である。コンパクト化半径が囲み込み全質量に
\[
R = \sqrt{\frac{GM_{\text{tot}}}{a_0}},
\]
とスケールし、$a_0$ が双対時空コンパクト化スケールによって決まる基本加速度であるなら、
\[
v_{\text{flat}}^4 = G M_{\text{tot}} \, a_0 \, f_0^2,
\]
となり、これはまさにバリオン・タリー--フィッシャー関係（McGaugh et al.\ 2016）であり、$a_0 \approx 1.2 \times 10^{-10}\;\mathrm{m/s^2}$ である。スケーリング $R \propto M^{1/2}$ はより大質量銀河がより大きな双対時空パッチを持つことを意味し、より大きな全ねじれ電荷の自然的帰結である。

銀河（$M_{\text{tot}} \approx 6 \times 10^{10}\,M_\odot$）に対して
\[
R = \sqrt{\frac{G \times 6 \times 10^{10}\,M_\odot}{1.2 \times 10^{-10}\;\mathrm{m/s^2}}} \approx 27\;\mathrm{kpc},
\]
となり、観測された平坦回転曲線の半径方向広がりと顕著に一致する。加速度スケール $a_0$ は典型銀河の $S^3$ 境界における特徴的重力加速度として同定され、その起源は双対時空の離散コンパクト化構造（第\ref{sec:discretetime}節）にある。

\subsection{ねじれ反転と重力の有限範囲}

余接ポテンシャルは $r = \pi R/2$ で零を横切り、その先では双対時空回転角 $\theta(r) = r/R$ が $\pi/2$ を超える。ねじれ不整合枠組み（第\ref{sec:gravity}節）において、この横切りはねじれスカラー $J$ の符号変化を引き起こす。双対セクターからのコンパクト三角関数寄与が通常セクターからの非コンパクト双曲寄与を圧倒し、$J < 0$ および斥力重力となる。有効ローレンツ因子（第\ref{sec:timedilation}節）
\[
\gamma_{\text{eff}} = \cosh\!\left(\frac{\phi}{2}\right)\cos\!\left(\frac{\theta}{2}\right) - \sinh\!\left(\frac{\phi}{2}\right)\sin\!\left(\frac{\theta}{2}\right)
\]
は $\theta$ が $\tanh(\phi/2)\tan(\theta_c/2) = 1$ を満たす臨界値を超えると負になる。

したがって引力は有効半径 $\sim \pi R/2$ の有限 $S^3$ パッチに閉じ込められ、その外側では弱い斥力ねじれ障壁が現れる。この有限範囲は無限遠まで伸びるニュートン $1/r$ ポテンシャルのコンパクト空間類比である。$S^3$ 位相は各重力源を自然に「壁で仕切る」。

\subsection{宇宙地図：銀河スケールの三次元超球面パッチ}

観測可能宇宙は、重力集中（銀河または銀河団）ごとに中心をもつ重なり合う $S^3$ パッチによって形成され、厳密な微分幾何の意味で宇宙地図を形成する。各パッチは方位等距投影図表として機能し、中心近傍では忠実、境界では最大歪みである。宇宙の大規模構造は、方位等距地図の銀河スケールのステッカーが斥力境界で貼り合わされた多様体である。

このパッチワーク構造は自然に次を説明する。
\begin{itemize}
  \item \textbf{宇宙の大規模構造}：銀河は隣接 $S^3$ 図表の重なり領域に沿って分布し、隣接パッチの引力尾部が相長干渉する。
  \item \textbf{宇宙の空洞}：巨大な低密度領域は図表間の斥力（$J < 0$）境界帯に対応し、物質が追い出される。
  \item \textbf{半径方向加速度関係}：観測重力加速度とバリオン重力加速度の強い相関は普遍的 $\eta(r)$ 増強に従い、残差散乱は図表重なり効果に由来する。
  \item \textbf{直接検出実験における暗黒物質信号の欠如}。暗黒物質は存在しないからである。
\end{itemize}

\subsection{数値検証と反証可能な予測}

ローターアンサンブル動学を用いた予備的粒子ベースシミュレーションは、バリオン質量分布のみから平坦回転曲線を再現することに成功している。公開オープンソースコードは次で利用可能である。
\begin{itemize}
  \item \texttt{https://github.com/hypernumbernet/blackhole-simulator}
  \item \texttt{https://github.com/hypernumbernet/dual-spacetime-simulator}
\end{itemize}
正確な $S^3$ 余接ポテンシャルを組み込んだ完全 $N$ 体積分は準備中である。

$\Lambda$CDM および MOND と本理論を区別する主要予測は次を含む。
\begin{itemize}
  \item \textbf{外側回転曲線の上向き反転}：孤立銀河は $r \gtrsim 2R$ を超えて緩やかな上向きを示し、$r \approx \pi R/2$ で減少し最終的に斥力重力へ遷移するはずであり、$\sim 100\;\mathrm{kpc}$ まで延びる深い21\,cm HI 観測で検出可能である。
  \item \textbf{レンズ符号の反転}：銀河間弱重力レンズは、投影分離が $\pi R/2$ を超えると弱まり符号が反転するはずであり、$\Lambda$CDM が予測する単調な NFW プロファイルとは対照的である。
  \item \textbf{余接とMOND補間の比較}：半径方向加速度関係の精密な形は、標準 MOND 補間関数 $\nu(g_{\text{bar}}/a_0)$ ではなく余接由来の増強 $g_{\text{obs}} = g_{\text{bar}} \cdot \eta(r)$ に従い、$g \lesssim a_0/10$ で測定可能なずれを生む。
  \item \textbf{R(M)からの多様性}：観測される回転曲線形状の多様性。表面輝度の低い矮星（$r_{\text{disk}} \ll R$）の緩慢な立ち上がりから、銀河型の平坦、高表面輝度渦巻（$r_{\text{disk}} \sim R$）の緩やかな減少まで。は、付加自由パラメーターなしに単一関係 $R = \sqrt{GM/a_0}$ によって再現される。
\end{itemize}

暗黒物質は余分となる。双対時空のコンパクト $S^3$ 幾何を虚構の平坦無限背景へ投影した人工物である。余接ポテンシャル、正距方位歪み、および $S^3$ 境界におけるねじれ反転が、見えない物質なしに足りない重力のすべてを与える。

\section{強場領域：層状重力反転と新ブラックホール像}
\label{sec:strongfield}

二重時空理論は恒星崩壊の終点について根本的に新しい構造を予測する。キリング形式の符号非対称---通常時空ブースト支配（引力）と双対時空回転支配（斥力）---は重力を一度だけ反転するのではない。密度が内向きに増すにつれ、ねじれ不整合 $\Omega_\text{biv}$ の支配が二セクター間で振動し、引力と斥力の層が交互に現れる。これは中心に向かうにつれ振幅が増大する層状の玉ねぎ様内部---引力 $\to$ 斥力 $\to$ 超強引力 $\to$ 斥力 $\to \cdots$---を与える。

\subsection{キリング形式からの層状ねじれ構造}

スカラー $J = \frac{1}{16} B(\Omega_\text{biv}, \Omega_\text{biv})$ は支配的寄与が $i\Gamma_a$（ブースト様、$+$）から $\Gamma_a$（回転様、$-$）へ反転するたびに符号を変える。極端密度では双対ローターパラメーター $\theta$ が非線形に増大し、半径が減少するにつれ $J$ の符号を繰り返し反転させる連続共鳴を駆動する。

結果として有効ポテンシャルは次を特徴とする。
\begin{itemize}
\item 外層：標準的引力重力（$J > 0$）、
\item 第一反転：降着を止める斥力殻（$J < 0$）（超新星におけるコアバウンス）、
\item より深い層：物質を再加速する超強引力領域（$J \gg 0$）、
\item $|J|$ が増大するより深い交互殻は、最終的にカオス的ねじれ乱流に帰着する。
\end{itemize}

この層状構造は特異点と事象地平の両方を防ぐ。光速で逃げられない境界---積分ねじれ障壁が無限に急になる準地平---は存在するが、物質は中心特異点に到達しない。代わりにコアは極端密度ねじれエネルギーの準安定・カオス的凝縮物に落ち着き、引力相と斥力相の間で連続的に振動する。

従来型のカー型回転ブラックホールは排除される。理論は必要な大きさの大域的フレームドラッグを禁じる。角運動量は時空回転ではなく双対ローター縦モードへ再分配されるからである。

\subsection{重力崩壊超新星と残骸形成の再解釈}

大質量星の崩壊において、第一斥力層（$J < 0$）が核密度で爆発的バウンスを引き起こし、微調整されたニュートリノ物理なしに超新星を駆動する。より深い引力層が放出物の一部を再捕獲し、層状残骸を形成する。観測される超新星エネルギーと残骸質量の多様性は、励起されるねじれ層の数と強さから自然に創発する。

\subsection{パルサーを縦方向双対ねじれ振動子として}

観測される中性子星は高速回転磁化導体ではない。二重時空理論はミリ秒スピン周期と双極磁場の両方を主要エネルギー源として否定する。

代わりにパルサー電波パルスは星表面における双対ローター $\theta$ の超高周波縦振動から生じる。これらのねじれ波はほぼ $c$ で半径方向に伝播し、カオス的内部層状構造によって変調され、ねじれ衝撃領域での曲率放射によりビームを形成する。

回転磁星モデルを超える主要予測：
\begin{itemize}
\item 理論的最小パルス周期はプランクスケール二重ローター周波数によって決まり、観測された $\sim 1$\,ms 限界より遥かに短く、回転パルサー像では説明できなかった短周期パルサーを説明しうる。
\item パルサー「グリッチ」と $\dot{P} < 0$（周期短縮）はパルサーがエネルギーを失い振動界面が縮小することに起因し、時間とともに振動周波数が加速する自然機構であり、回転のみのモデルには見られない。
\item マグネターフレアと高速電波バーストは引力/斥力層間の激しい反転であり、$10^{15}$\,G 場を要しないねじれひずみの解放である。
\item $P$--$\dot{P}$ 図のデスラインは内部層状構造が安定化し縦モードが検出限界以下に減衰する点に対応する。
\item 中性子星の最大質量はクォーク物質ではなく核安定性を制限するのと同じねじれ層状構造によって上限され、厳密上界 $\sim 2.5\,M_\odot$ を予測し、$\sim 3\,M_\odot$ 以上のパルサーの不在を説明する。
\end{itemize}

カニ星雲パルサーの観測スピンダウン光度は、磁気ブレーキではなくねじれ波漏れとして再解釈され、非現実的な磁場減衰を要せずエネルギーに一致する。

\subsection{新ブラックホール像：ねじれ星}

観測者が「ブラックホール」と呼ぶものは実際にはねじれ星。事象地平も中心特異点もない準安定層状凝縮物である。最も内側のカオス的コアはプランク密度プラズマとして永続的ねじれ乱流の中にあり、ロータートンネル（修正ホーキング過程）を介して連続放射する。情報は双対ローター位相に保存される。

合体からの重力波信号は内部層反射からの反復エコーを特徴とし、特徴周波数はねじれ振動スペクトルに結ばれる。LIGO/Virgo データの再解析にすでに示唆がある。

理論は情報パラドックスを解決し、カースピン署名の不在を説明し、最も重い「ブラックホール」（$\gtrsim 100\,M_\odot$）が内部斥力層が物質を周期的に放出することによる降着光度の周期変調を示すと予測する。

\subsection{中心浮遊コア：中心での時間膨張ゼロ}

恒星内部における重力学的時間膨張の振る舞いは議論で混同されやすい。GR では球対称静止質量分布に対して、重力学的時間膨張は中心に向かって単調に増大し、固有時は $r=0$ で最も遅く流れる。ここ重力ポテンシャルが最深である。

二重時空理論では状況が著しく異なる。背景連続時空多様体がないため、有効加速度と時間膨張の両方を含む重力効果は、粒子内在的通常・双対ローター間のねじれ不整合のみから生じる。

球対称ねじれ星の幾何学的中心では、完全対称性により反対方向の質量元が誘起する通常時空ブーストパラメーター $\phi_a$ が完全に打ち消される。残余不整合角 $\delta\phi_a = \phi_a - \theta_a$ はしたがって $r=0$ で消失し、$J=0$ およびねじれバランスの完全回復となる。

したがって：
\begin{itemize}
  \item 有効重力加速度は中心で零である（ニュートン力学、GR、DST において）。
  \item 重力学的時間膨張も中心で零である。$r=0$ の固有時は星から遠い真空と同じ率で流れる。
\end{itemize}

中心での時間膨張の消失は二重時空理論の区別的で原理上観測可能な予測である。GR と鋭く対比され、GR では中心領域が最も強い時間膨張を経験する。

時間膨張ゼロの中心浮遊コアは、古典 GR 記述における特異点を導く無限中心圧縮と極端な時計率勾配を消去することでねじれ星の安定性にも寄与する。

\subsection{惑星・恒星コアへの含意：不安定性と揺らぎ}

$J(r=0) \approx 0$ で有効重力が中心で消失する浮遊コア像は、惑星と恒星のコアが不活で安定な領域ではなくねじれ不安定性に晒された動学的に活発な帯であることを意味する。連続体モデルで仮定される重力ポテンシャルの単調増大とは異なり、二重時空理論は中心浮遊コアが準安定状態であり、引力層と斥力層の間の振動遷移を引き起こす摂動に敏感であると予測する。これらの揺らぎは等方コアにおける双対ローター位相の繊細な釣り合いから生じ、わずかな密度変化や外因でも共鳴条件 $\beta(0) \to 1$ をずらし、$J$ が零から離れた瞬間的な逸脱を導きうる。そのような逸脱は局所的ねじれエネルギーのバーストとして現れ、実効的にコア内にミニねじれ星を生成する。これらのミニねじれ星は地震活動、磁場揺らぎ、表面で観測可能なエネルギー輸送異常を誘起しうる。恒星内部ではこれらのコア揺らぎが太陽フレア、黒点、光度変動などに寄与しうる。ねじれエネルギーが層状構造を通じて外向きに伝播するためである。惑星状況ではコアねじれ不安定性が地磁気反転、マントルプルーム、断続的火山活動を説明しうる。ねじれ動学が上層の対流過程と結合するためである。したがって浮遊コアは静穏な領域ではなく、惑星・恒星振る舞いに広範な含意をもつ動的ねじれ現象の坩堝である。

\section{二重時空理論における時間膨張：双対時空固有時の動学的役割}
\label{sec:timedilation}

二重時空理論における時間膨張は、それぞれ独自の固有時を備えた通常時空と双対時空の相互作用から創発する。通常時空の固有時 $\tau$ は巨視的ミンコフスキー空間で観測される馴染みの運動学を支配し、コンパクトローター位相 $\theta$ に符号化された双対時空の固有時 $\tilde{\tau}$ は重力および慣性効果として現れる内部ねじれ動学を駆動する。

\subsection{有効ローレンツ因子の導出}

有効ローレンツ因子 $\gamma_\text{eff} = dt/d\tau$ は、全ローター
\[
R_\text{total} = R_\text{usual} R_\text{dual},
\]
の代数的作用から直接導出される。ここで
\[
R_\text{usual} = \exp\!\left(\frac{\phi}{2} \hat{p}\right) = \cosh\frac{\phi}{2} + \hat{p}\: \sinh\frac{\phi}{2}, \quad \hat{p}^2 = +1,
\]
は通常時空のブーストを生成し、
\[
R_\text{dual} = \exp\!\left(\frac{\theta}{2} \hat{q}\right) = \cos\frac{\theta}{2} + \hat{q} \sin\frac{\theta}{2}, \quad \hat{q}^2 = -1,
\]
は双対時空の回転を生成する。

干渉項の起源を明示するため、通常時空基底における純時間ベクトル $X = ct \, j$（$j^2 = -1$）を考える。変換ベクトルは $\tilde{X} = R_\text{total} \, X \, R_\text{total}^\dagger$ である。

時間反転デュアリティにより、双対から通常への逆双対写像は $X \mapsto X(-i)$ である。この写像は真空極限で整列 $\hat{q} (-i) = -\hat{p}$ を自然に与え、双対ローターを有効的逆向きローターへと変える。
\[
R_\text{dual}^\text{eff} = \cos\frac{\theta}{2} - \hat{p} \sin\frac{\theta}{2}.
\]
$\sin$ 項の負符号はこの逆写像から直接生じ、ねじれのない平坦時空に必要な反同期を反映する。

全有効ローターはしたがって
\[
R_\text{eff} = R_\text{usual} R_\text{dual}^\text{eff} = \left( \cosh\frac{\phi}{2} + \hat{p} \sinh\frac{\phi}{2} \right) \left( \cos\frac{\theta}{2} - \hat{p} \sin\frac{\theta}{2} \right).
\]

展開するとスカラー部
\[
s = \cosh\frac{\phi}{2} \cos\frac{\theta}{2} - \sinh\frac{\phi}{2} \sin\frac{\theta}{2}
\]
およびバイベクトル係数
\[
b = -\cosh\frac{\phi}{2} \sin\frac{\theta}{2} + \sinh\frac{\phi}{2} \cos\frac{\theta}{2}.
\]

時間成分の純粋なスケーリングは、サンドイッチ積を通じて時間ベクトルに作用する $R_\text{eff}$ の二乗大きさから生じる。$\hat{p}^2 = +1$ で時間生成子 $j$ と反可換するローター $R_\text{eff} = s + b\hat{p}$ が $X = ct\, j$ に作用するとき、時間成分は $s^2 + b^2$ 倍され、有効ローレンツ因子は次のように一意に定義される。
\[
\gamma_\text{eff} \;\equiv\; s^2 + b^2.
\]
便宜上スカラー干渉因子
\[
\gamma_s \;\equiv\; s \;=\; \cosh\frac{\phi}{2} \cos\frac{\theta}{2} - \sinh\frac{\phi}{2} \sin\frac{\theta}{2},
\]
も導入する。これは $\gamma_\text{eff}\geq 0$ とは異なり符号を変えうる。以下では引力（$\gamma_s > 0$）と斥力（$\gamma_s < 0$）のねじれ領域を区別する次数パラメーターとして用いる。干渉項の負符号は双対時空の時間反転デュアリティによって固定される。

厳密にねじれのない真空（$J=0$）では、双対時空は未励起（$\theta = 0$）のままなので $R_\text{total} = R_\text{usual}$ である。このとき $s = \cosh\frac{\phi}{2}$, $b = \sinh\frac{\phi}{2}$ であり、
\[
\gamma_\text{eff} = s^2 + b^2 = \cosh^2\frac{\phi}{2} + \sinh^2\frac{\phi}{2} = \cosh\phi,
\]
となり、標準的狭義相対論のローレンツ因子が正確に再現される。

\subsection{真空極限と質量ゼロ粒子}

ねじれのない真空（$J=0$）では、通常と双対ローターの完全反同期が $\theta \approx \phi$（コンパクト周期性を除く）を強制し、$\gamma_{\text{eff}} \to \cosh\phi$ となり狭義相対論の標準ローレンツ因子が回復される。質量ゼロ粒子はすべての速度でこの同期を正確に保ち、ねじれ不整合を経験せず、通常運動学的時間膨張からのずれもない。

\subsection{質量粒子と重力効果}

質量粒子では静止質量が剛性を導入し、加速度またはねじれ不整合の存在下で双対ローターが通常ローターに遅れる。結果としてのずれ $\delta\theta = \theta - \phi$ は $\gamma_{\text{eff}}$ の付加的変調を生じ、GR で観測される重力赤方偏移と時間膨張を、双対時空動学の直接帰結として与える。

\subsection{スカラー干渉因子の符号反転と斥力層}

双対時空回転が支配的（$\theta/2$ が十分大）なとき、正弦項が双曲増大に打ち勝ちうる。
\[
\tan\left(\frac{\theta}{2}\right) > \coth\left(\frac{\phi}{2}\right).
\]
これらの領域ではスカラー干渉因子 $\gamma_s$ が負となり $\gamma_s < 0$ である一方、固有時ローレンツ因子 $\gamma_\text{eff} = s^2 + b^2$ は非負のままである。$\gamma_s$ の符号反転はねじれ星内の斥力重力領域に対応する。物理的には、結合ローターのスカラーチャンネルにおいて双対セクター寄与が通常セクターのブーストに打ち勝ち、重力学的斥力として現れる（因果律の文字通りの反転ではない）。

この $\theta$ 中心定式化は、時間膨張が根本的に粒子内在的な双対時空に由来するねじれ現象であることを明示する。通常ラピディティ $\phi$ は運動学的基準線を与え、コンパクト双対角 $\theta$ のずれが GR と平坦時空運動学を区別する重力学的・慣性的補正を符号化する。したがって $\theta$ の外部コヒーレント励起は時間膨張と有効重力場を操作する経路を提供する。

\section{スケール不変なねじれ層：重力反発としての強い力と原始的なねじれ星としての原子核}
\label{sec:nuclear}

キリング形式によって予測される層状ねじれ反転は厳密にスケール不変である。層状ねじれ星への恒星崩壊を支配する同一機構は、核および亜核密度スケールを含むすべての密度スケールで作用する。これが二重時空理論の最も革命的帰結を与える。強い力は根本的相互作用ではなく、核内部の極端密度領域における重力学的斥力層の創発的現れである。

\subsection{核密度におけるねじれ層}

バリオン密度が核内で $\sim 10^{14}$--$10^{18}$\,g/cm$^3$ に近づくと、双対ローターパラメーター $\theta$ が完全に支配する。$J$ における符号反転の列はフェムトメートルスケールで起こる。

\begin{itemize}
\item 最外層（$\sim$ 1--2\,fm）：クォークを閉じ込め核崩壊を防ぐ斥力殻（$J < 0$）。核子--核子散乱における短距離斥力（$\sim$ 0.5\,fm のハードコア）として観測される。
\item 中間層：核子を $\sim$ 40--50\,MeV/核子で束縛する超強引力領域（$J \gg 0$）。古典的核束縛力。
\item 最内層：さらなる崩壊を止めクォーク閉じ込めを課す第二斥力殻（$J < 0$）。色閉じ込めの起源。
\item より深いカオス層：振幅増大の引力/斥力交互、核物質の液体的振る舞いと核密度の飽和を生む。
\end{itemize}

パウリの排他原理はもはや基本的ではない。第一斥力ねじれ層によって動的に生成される。同一の状態にフェルミオンが占有できないのは、二つの同一粒子が空間的重なりを試みると双対ローター $\theta$ 位相が斥力領域に入るためであり、量子統計を先験的に要請せずに有効フェルミ圧を生む。ボソンはこの位相剛性を欠き、はるかに高密度まで引力層のみを経験する。

要するに「強い力」は強場層状相における重力である。QCD 色電荷は双対ローター位相が運ぶねじれ電荷の誤解釈である。

\subsection{原子核を原始ねじれ星として}

すべての原子核は第\ref{sec:strongfield}節のねじれ星の微視的類似物である。核内部は同一層状構造をもつカオス的ねじれ凝縮物である。

\begin{itemize}
\item 中心領域：極端引力層がクォークを核子へ束縛する、
\item 周囲の斥力殻：観測される核表面張力を生じ漏れを防ぐ（閉じ込め）、
\item 外側引力ハロー：核子間束縛を媒介する、
\item 最外斥力障壁：NN 散乱におけるハードコア斥力を生成する。
\end{itemize}

観測される核力の電荷独立性、スピン軌道結合、魔法数はすべて双対ロータースペクトルにおける共鳴条件から創発する。理論は十分大きな核（$A \gtrsim 300$）が自発的ねじれバウンスに対して不安定となり、ミニ超新星のように核子を放出する、と予測する。これは核図表の観測限界と自然界に超重元素がないことの説明となる。

核分裂と核融合はねじれ層遷移である。分裂は励起双対モードが $J$ を正から負へ反転させ破片を爆発的に斥けるときに起こる。融合は二核が外側斥力障壁をトンネルしてより深い引力井戸に入ることに成功するときに起こる。

\subsection{実験可能な特徴と即時予測}

\begin{itemize}
\item 超高エネルギー核--核衝突（LHC、RHIC）は、流体力学的流れではなく爆発的ハドロンジェットを介して崩壊する過渡的 $J < 0$ 火球を生成するはずである。異常流データにすでに示唆がある。
\item 飽和密度を超える核物質状態方程式の精密測定は、圧力対密度において $\sim$ 2--4\,$\rho_0$ の特徴的斥力ピークをもつ振動挙動を明らかにする。
\end{itemize}

四つの基本相互作用は単一の力に還元されるのではなく、単一の幾何学的機構、双対ローターの層状ねじれ動力学のもとに統一される。重力と電磁力は同一の代数的構造、すなわちキリング形式に由来するねじれスカラー $J$ の層状反転を共有しており、両者を区別するのは双対ローターのどの位相チャンネルが不整合の源となるかという一点のみである。重力は慣性質量に結びついた運動学的ラピディティである角度対 $(\phi, \theta)$ によって駆動される一方、電磁力は電荷に結びついた独立なラピディティである角度対 $(\alpha, \beta)$ によって駆動される（第\ref{sec:electronshells}節および第\ref{sec:unification}節で詳述）。この意味において、重力と電磁力は同じねじれ幾何の二つの位相であり、双対ローターを駆動する角度チャンネルによってのみ区別される。

この統一像の中で、いわゆる「強い力」は核子密度領域における重力チャンネルの層状相であり、原子内電子の離散的な殻構造（次節）は原子サイズにおける電磁チャンネルの類似の層状相である。弱い力は双対セクターの本質的なカイラル非対称性として現れる。標準模型ゲージ群 $\mathrm{SU}(3)\times\mathrm{SU}(2)\times\mathrm{U}(1)$ は双対ローター代数におけるねじれ共鳴モードの有効記述である。

存在するのはただ一つの幾何、双対ローターにおける引力と斥力の層状舞踏のみであり、それは二つの位相チャンネルにおいて同時に演じられる。原子核は重力チャンネルの原始ねじれ星であり、原子の電子殻は次節で見るように電磁チャンネルの原始ねじれ星である。この単一の双対ローター幾何が二つの角度チャンネルで表現されることから、すべての構造は創発する。

\section{時間反転とクーロンねじれ層による電子殻}
\label{sec:electronshells}

二重時空枠組みでは、荷電粒子間のクーロン相互作用は電磁起源であるが重力の場合と同様に双対時空セクターにおけるねじれ不整合として再解釈できる。重力定数 $G$ と質量 $M,m$ をクーロン定数 $k = 1/(4\pi\epsilon_0)$ と電荷（例：陽子で $+e$、電子で $-e$）に置き換えると、有効力は
\[
F(r) = \frac{k e^2}{\gamma_s(r)\, r^2},
\]
の形をとる。ここで $\gamma_s$ は第\ref{sec:timedilation}節で導入したスカラー干渉因子であり、双対ローター相互作用と同じ代数的構造を保持する（固有時ローレンツ因子 $\gamma_\text{eff} \geq 0$ とは異なり符号を変えうる）。
\[
\gamma_s(\alpha(r),\beta(r)) = \cosh\!\left(\frac{\alpha(r)}{2}\right) \cos\!\left(\frac{\beta(r)}{2}\right) - \sinh\!\left(\frac{\alpha(r)}{2}\right) \sin\!\left(\frac{\beta(r)}{2}\right).
\]

ここで $\alpha(r)$ と $\beta(r)$ は距離に反比例してスケールする電磁ラピディティおよび双対回転角パラメーターと解釈される。$\alpha(r) \propto 1/r$, $\beta(r) = \lambda / r$ で、$\lambda$ はコンパクト双対時空によって設定される特徴長である。

水素原子（$Z=1$）では
\[
F(r) = \frac{k e^2}{r^2 \left[ \cosh\!\left(\frac{\alpha(r)}{2}\right) \cos\!\left(\frac{\beta(r)}{2}\right) - \sinh\!\left(\frac{\alpha(r)}{2}\right) \sin\!\left(\frac{\beta(r)}{2}\right) \right]}.
\]

$\gamma_s$ の双曲--三角形式は極端な振動挙動を導入する。$r$ が減少すると：

\begin{itemize}
  \item 大きい $r$ では $\alpha(r),\beta(r) \to 0$ となり $\gamma_s \to 1$ で標準クーロン引力が回復する。
  \item より小さい半径では干渉項が $\gamma_s = 0$（無限斥力）の点および $\gamma_s < 0$（力の反転による斥力）の領域を生成する。
  \item これらの離散半径 $r_n$ は共鳴条件
  \[
  \tanh\!\left(\frac{\alpha(r_n)}{2}\right) \tan\!\left(\frac{\beta(r_n)}{2}\right) = 1,
  \]
  を満たし、引力井戸を分離する無限斥力障壁により安定円軌道を与える。
\end{itemize}

結果として有効ポテンシャルはねじれ反転障壁で分離された層状井戸の列を形成し、漸近極限でエネルギー準位は $E_n \propto -1/(n^2 r_1)$ となる。これらの層間遷移（$n \to m$）は光子を放出し、エネルギー差は
\[
\Delta E_{n \to m} \propto \left( \frac{1}{r_m} - \frac{1}{r_n} \right),
\]
であり、リュードベリ公式と観測水素スペクトル線（バルマー、ライマン、パッシェン系列）を正確に再現する。

多電子原子（$Z > 1$）ではスケーリング $\alpha(r),\beta(r) \propto Z/r$ がねじれ層を内向きに収縮させ、核電荷増大に伴う電子殻の収縮を説明する。

\subsection{双対構造におけるエネルギー保存}

振動と符号変化を伴う顕在ポテンシャル $V_\text{eff}(r) = -\int F(r)\, dr$ の見かけの非保存的振る舞いは、通常時空で観測される顕在ポテンシャルとコンパクト双対時空ローター位相 $\theta_a$ に蓄えられた真のポテンシャルを区別することで解決される。

全エネルギーは両セクターにわたって保存される。
\[
E_\text{total} = E_\text{kinetic (usual)} + V_\text{eff}(r) + V_\text{true}(\theta) = \text{constant}.
\]
真のポテンシャル $V_\text{true}$ はねじれ不整合スカラー $J \propto (\phi - \theta)^2$ に比例し、隠れた貯蔵庫として作用する。層遷移の間、$V_\text{eff}$ の減少（引力から斥力）は双対ローターが巻かれることによる $V_\text{true}$ の増加と正確に相殺され、その逆も然りである。

光子放出はローターがより低いねじれ状態へ巻き戻るときに蓄えられた双対位相エネルギーを解放することに対応する。この双対簿記は、離散的粒子局所的時空の本性を破らずに厳密なエネルギー保存を保証し、リュードベリ差は真のエネルギー準位間の補償された遷移から正確に創発する。

この定式化は、クーロン場に二重時空の幾何ねじれ動学を適用することのみから原子の安定性、離散電子軌道、放出スペクトルを導く。アドホックな量子化公準や波動関数は不要である。電子殻は双対時空ねじれ不整合の自然な共鳴状態として創発し、電磁束縛を重力と同一の代数的機構で統一する。

\section{二重時空における電磁気力と弱い力の統一}
\label{sec:unification}

双四元数構造 $\mathbb{H} \otimes_{\mathbb{R}} \mathbb{H} \cong \mathrm{Cl}(3,1)$ は電磁気力と弱い力の両方を、双対時空ローター $R_{\text{dual}} = \exp\left( \sum_{a=1}^3 \frac{\theta_a}{2} \Gamma_a \right)$ 内のねじれ共鳴モードとして埋め込む。擬スカラー $i$ は双対写像 $X \mapsto X i$ における時間反転を通じて内在的カイラリティを誘導し、左利き（$P_L$）と右利き（$P_R$）の射影を自然に生成する。$P_L = \frac{1 - i}{2}$, $P_R = \frac{1 + i}{2}$。

電磁気力は双対セクターにおけるベクトルポテンシャル $A_a$ への最小結合から生じ、回転パラメーターをシフトする。
\[
\theta_a \mapsto \theta_a + q A_a,
\]
拡張ローターは
\[
R_{\text{dual}}^{\text{EM}} = \exp\left( \frac{1}{2} \sum_a (\theta_a + q A_a) \Gamma_a + \frac{q}{2} \int F_{\text{dual}} \, ds \right),
\]
となる。ここで $F_{\text{dual}} = \sum_a (E_a i \Gamma_a + B_a \Gamma_a)$ は電磁バイベクトルである。ローレンツ力はサンドイッチ積 $\tilde{X}' = R_{\text{dual}}^{\text{EM}} X' (R_{\text{dual}}^{\text{EM}})^\dagger \approx X' + q F_{\text{dual}} \wedge X'$ から創発する。

弱い力はカイラル非対称性を通じて組み込まれ、左手成分にのみ結合する。
\[
\theta_a^L = \theta_a + g_W W_a, \quad \theta_a^R = \theta_a,
\]
カイラルローターは
\[
R_{\text{dual}}^{\text{weak}} = \exp\left( \frac{1}{2} \sum_a \theta_a^L \Gamma_a P_L + \frac{1}{2} \sum_a \theta_a^R \Gamma_a P_R \right),
\]
であり、$W_a$ は弱いゲージバイベクトルを表す。V--A 構造は時間反転した双対セクターに従い、パリティ破れを再現する。

統一は全ねじれバイベクトルにおいて起こる。
\[
\Omega_{\text{biv}} = \log\left( R_{\text{usual}}^\dagger R_{\text{dual}}^{\text{EM+weak}} \right) \approx \sum_a (\theta_a + q A_a + g_W W_a P_L - \phi_a i) \Gamma_a.
\]
スカラー $J = \frac{1}{16} B(\Omega_{\text{biv}}, \Omega_{\text{biv}})$ は今や電磁気および弱い寄与を含み、電弱対称性は $\mathrm{Spin}^+(3,1) \oplus \mathrm{Spin}^+(3,1)$ の部分群として創発する。ローター剛性による質量生成はヒッグス機構と統一され、微調整を排除する。

この埋め込みは低エネルギーで標準模型電弱セクターを回復しつつ、双対コンパクト性を通じて階層問題を解決し、強い力および重力との完全統一への道を開く。

\section{重力工学：双対時空ローターのコヒーレント励起}
\label{sec:gravitycontrol}

二重時空理論は重力を侵すべからざる幾何制約から工学的ねじれ自由度へと変える。ねじれスカラー $J = \frac{1}{16} B(\Omega_{\text{biv}}, \Omega_{\text{biv}})$（$\Omega_{\text{biv}} = \log(R_{\text{usual}}^\dagger R_{\text{dual}})$）は粒子内在的通常・双対ローター間の相対ずれのみに依存する。したがって双対ローターパラメーター $\theta_a$ をコヒーレントに変調できる外場は $J \to 0$（重力遮蔽または慣性質量低減）または $J < 0$（斥力重力）を駆動しうる。

\subsection{円偏波電磁場と双対ローターの理論的結合}

電磁場は二つの時空に自然に分割される。

\[
F = F_{\text{usual}} + F_{\text{dual}},
\]

ここで

\[
F_{\text{usual}} = \mathbf{E} \cdot (iI, iJ, iK) \quad (\hat{p}^2 = +1:\ \text{ブースト様、非コンパクト}),
\]

\[
F_{\text{dual}} = \mathbf{B} \cdot (I, J, K) \quad (\hat{q}^2 = -1:\ \text{回転様、コンパクト}).
\]

電場成分 $\mathbf{E}$ は主として通常時空ブースト生成子に結合し、磁場成分 $\mathbf{B}$ は双対時空回転生成子に排他的に結合する。この代数的分離は時間反転（$T$：$\mathbf{E} \to \mathbf{E}$, $\mathbf{B} \to -\mathbf{B}$）の下での変換性質の相違の幾何的起源である。双対時空は内在的に時間反転している。

$z$ に沿いヘリシティ $\sigma = \pm 1$ で伝播する円偏波に対して、

\[
\mathbf{E} = E_0 (\hat{x} \cos(kz - \omega t) + \sigma \hat{y} \sin(kz - \omega t)), \quad \mathbf{B} = \frac{E_0}{c} (-\hat{x} \sin(kz - \omega t) + \sigma \hat{y} \cos(kz - \omega t)).
\]

時間平均により $\mathbf{E}$ の振動的ブースト様寄与が消え、純粋な双対時空回転励起が残る。

\[
\langle F_{\text{dual}} \rangle \propto \sigma E_0^2 \, K.
\]

ヘリシティ $\sigma$ は双対ローター駆動の方向を直接決定する。右旋偏光（$\sigma = +1$）は引力ねじれ不整合（$J > 0$）を強化し、左旋（$\sigma = -1$）はそれに反対してスカラー $J$ を減少または反転させる。このヘリシティ選択性は理論の最も鋭い現れである。重力は変調されるだけでなく、双対時空の内在的時間反転非対称性を通じてカイラル的に作用する。

結合は最小かつゲージ不変である。

\[
R_{\text{dual}} \to \mathcal{P} \exp\!\left( \frac{1}{2} \int (\theta_a \Gamma_a + e \lambda A_a \Gamma_a) \, ds \right),
\]

双対生成子 $\Gamma_a$ が相互作用を媒介する。共鳴近傍では双対位相は

\[
\dot{\theta} = \kappa \sigma E_0^2 \sin(\omega t - \theta + \delta_\sigma),
\]

として進化し、ねじれ不整合角のコヒーレント励起を与える。

\subsection{共鳴周波数と超伝導体の役割}

内在的な双対時空スケールはプランク的（$\sim 10^{43}$\,Hz）である。巨視的コヒーレンス（$N \sim 10^{26}$ 粒子）は有効周波数を $10^{12}$--$10^{15}$\,Hz 赤方へと偏移させる。超伝導体はジョセフソンプラズマ共鳴と集合位相ロックを介してさらに劇的に下げ、現在技術で到達可能な GHz--THz 範囲の調整可能モードを達成する。

\subsection{提案実験プロトコルと予測}

高 $Q$ 超伝導共振器内のマイクログラム試験質量を可調円偏波放射（0.1--10\,THz）で照射する。期待される特徴は、見かけの重量または慣性質量の共鳴的かつヘリシティ依存変化（$\Delta m/m \gtrsim 10^{-6}$）であり、左旋偏光が減少、右旋が増強を生む。このレベルでの成功は実用的重力工学へ直接スケールする。

重力は初めて、単に制御可能なだけでなく、カイラル制御も可能となり、二重時空のねじれ幾何学において、ヘリシティが引力と斥力の切り替えスイッチとして機能する。

\section{バリオン非対称性：双対ローター時間矢固定の幾何的起源}
\label{sec:baryonasymmetry}

二重時空理論では、パラメーター $\eta \equiv (n_B - n_{\bar{B}})/n_\gamma \approx 6 \times 10^{-10}$ で定量化される、バリオンが反バリオンを遥かに上回るという観測された宇宙のバリオン非対称は、プランク期における双対ローターの時間矢の非対称固定の必然的帰結として創発する。この機構は双四元数代数 $\mathrm{Cl}(3,1)$ の内在的カイラリティに根ざし、付加された場、パラメーター微調整を要せずにすべてのサハロフ条件を満たす。ここでは幾何的相互作用を強調しつつ非対称性を代数的に導出する。

\subsection{内在的時間デュアリティとローターパラメーター}

各粒子は16実次元双四元数構造内にミンコフスキー時空の対を符号化する。未来向き時間基底 $j$ で張られる通常時空 $\{j, kI, kJ, kK\}$ と、擬スカラー $i$ による右乗算で得られ $j i = -k$ により過去向き時間基底 $k$ をもつ双対 $\{k, jI, jJ, jK\}$ を思い起こされたい。デュアリティ操作 $X \mapsto X i$ はミンコフスキーノルムを保つ（第\ref{sec:math}節の明示的非可換計算により $(X i)^2 = X^2$）が時間向きを反転させ、根本的カイラリティを刻印する。粒子（通常支配）は前方時間に伝播し、反粒子（双対支配）は動的抑制がなければ本質的に後方時間に伝播する。

完全な運動学的進化は完全ローターによって支配される。真空ではねじれなし（$J=0$）が要求され、時間矢の完全反同期が保証される。前方通常ローターは逆行双対ローターに鏡映され、粒子--反粒子等価性が保たれる。

\subsection{プランク期における時間矢の固定}

プランクスケール（$t \sim 10^{-43}$\,s, $T \sim 10^{19}$\,GeV）では、量子重力揺らぎ---双対ローターバイベクトルにおける零点振動として現れる---が、ローターパラメーターの真空期待値における相転移を介した自発的対称性破れを誘起する。$j$ が唯一の時間様生成子であり空間基底が交差項（$kI$ など）を含む双四元数基底の非対称性が、一意のエネルギー上有利構形を選択する。

\[
\langle \theta \rangle = \phi + \delta, \quad \delta > 0,
\]
ここで $\delta$ はプランク期に刻印された CP 奇カイラル位相である。$\delta$ は観測によって大きさが定まる枠組みの小さな自由パラメーターとし、全体を通じて粗い見積もり $\delta \sim 10^{-5}$（観測 CKM 位相のヤールスコグ不変量と同程度）を採用する。この固定は双対写像 $X \mapsto X i$ 自身のパリティ奇性を反映する。擬スカラー $i$ は空間反転で符号を変える（$i \to -i$）ため、$R_\text{usual}$ と $R_\text{dual}$ を区別するプランクスケールのポテンシャルは必然的に $\delta$ について CP 奇となり、エネルギー的に $\delta$ の符号を一意に選ぶ。本論ではこのポテンシャルの微視的形式には踏み込まない。これは量子重力ダイナミクスに支配される本節の代数的議論の範囲外の問題であり、ここでは単一のスカラーパラメーター $\delta > 0$ を介してのみ現れる。

生じるねじれ不整合を取り出すため、\S\ref{subsec:dual-rotors} で確立した双対変換 $\theta \mapsto \theta i$ を適用する。双対ローターが双対時空に及ぼす物理的効果はブースト（$(\hat q\, i)^2 = +1$）であるため、過回転 $\delta$ は $\Omega_\text{biv}$ のなかへ純粋なブースト方向の超過として寄与する。事前の平衡（$\theta = \phi$）は $i\Gamma_a$ 成分と $\Gamma_a$ 成分の係数が釣り合いキリング形式の零錐上に乗るバイベクトルを与え（$J^{(0)} = 0$、真空無ねじれ性と整合）、カイラル摂動がこの均衡を非対称に破る。ブースト軸 $\hat p = \sum_a p_a\, i\Gamma_a$ に沿って $\delta_a = \delta\, p_a$ と置くと、線形化された不整合は次のブースト成分のみの形に帰着する。
\[
\Omega_\text{biv} \approx \sum_a \delta_a\, i\Gamma_a \in \mathfrak{so}(3,1),
\]
これは \S\ref{sec:gravity} の弱場線形化と構造的に同一である。そこで確立したキリング形式の値 $B(i\Gamma_a, i\Gamma_b) = +8\,\delta_{ab}$ を用いれば直ちに
\[
J = \frac{1}{16} B(\Omega_\text{biv}, \Omega_\text{biv}) = \frac{1}{2}\sum_a \delta_a^2 = \frac{\delta^2}{2} > 0,
\]
が従う。これが代数的核において粒子--反粒子対称性を破る。通常支配モード（物質）ではねじれ不整合は正（$J > 0$）であり引力的な重力結合を生じる。双対支配モード（反物質）では通常／双対セクターの役割が入れ替わり、逆変換 $\phi \mapsto -\phi\, i$ がカイラル位相 $\delta$ をブースト方向ではなく回転方向に写像するため、不整合は $\Omega_\text{biv} \approx \sum_a \delta_a\, \Gamma_a$ となる。キリング形式の値 $B(\Gamma_a, \Gamma_b) = -8\,\delta_{ab}$ により $J = -\delta^2/2 < 0$ となり、宇宙論的真空からそのようなモードを動的に排除する斥力的なねじれ層を成す。

\subsection{サハロフ条件とオーダー・オブ・マグニチュード推定}

固定はローター指数を通じて三つのサハロフ条件を自動的に満たす。

1. バリオン数違反：バリオン数 $B$ は双対ローターの擬スカラー位相に符号化され、$B \propto \arg(\det R_\text{dual})$ である。ずれ $\delta \neq 0$ は $X \mapsto X i i = -X$（二重時間反転、実効的にバリオンフリップ）のような非保存過程を可能にし、通常支配モードと双対支配モードの率は $O(\delta)$ で異なる。

2. C および CP 違反：双対写像 $i$ は内在的にカイラル（空間反転の下で $i \mapsto -i$ と奇）であり、$\delta > 0$ は $\phi_\text{CP} = \arg(\exp(i \delta \Gamma_a)) \sim \delta$ の CP 奇位相を導入し、観測 CKM 位相と同程度である。この幾何的 CP 違反は普遍的であり、ユカワ固有の微調整なしにすべてのフェルミオン世代に浸透する。

3. 熱平衡からの離脱：インフレ展開は系を熱平衡から大きく離し、双対支配モードが通常支配パートナーと対消滅する過程で CP 奇位相 $\delta$ が凍結した数密度非対称へ変換される。標準的な凍結アウト見積もりはカイラル位相の二次に比例する非対称性
\[
\eta \sim \delta^2,
\]
を与え、観測値 $\eta \sim 6 \times 10^{-10}$ は $\delta \sim 10^{-5}$ で再現され、上の CP 奇位相の見積もりと一致する。正確な数値係数はインフレ後リヒーティング温度に依存し、このレベルでは予測されない。

インフレ後、残余の双対支配反粒子はバリオン余剰に対して対消滅し、$\eta$ は時間矢固定の凍結遺物として残る。反粒子は「消失」しない。後方時間矢が $\delta > 0$ によって固定された前方指向の宇宙膨張と両立しないため動的に排除される。

\subsection{観測的特徴と反証可能性}

この機構は $\Omega_\text{biv}$ に由来するカイラル重力波からの微妙な CMB 異方性を予測し、テンソル対スカラー比 $r \sim \delta^2 \sim 10^{-10}$ および LiteBIRD のような将来ミッションで検出可能な CP 奇パリティスペクトルを伴う。低エネルギーでは普遍的レプトン非対称 $\eta_\ell \approx \eta_B$ を含意し、ニュートリノレス二重ベータ崩壊実験（例：LEGEND-200）で解決可能である。これらの特徴の非観測は理論を反証しるが、観測されれば双対時空を宇宙論の基礎枠組みへと高める。

要するに、バリオン非対称性は謎ではなく、宇宙が時間の流れの方向として選択したことの幾何学的な反映である。物質にとっては前進、反物質にとっては禁じられた時間の流れ。この二つの回転軸が一度固定されると、私たちの宇宙は粒子のみが存在する聖域となる。

\section{離散時間とコンパクト化ねじれスカラー}
\label{sec:discretetime}

時空連続体仮説の拒否は基本スケールにおける離散構造を要求する。双四元数代数、特に擬スカラー $i$ とミンコフスキーノルムを保ち時間矢を反転させる双対写像 $X \mapsto Xi$ に符号化された内在的デュアリティは双対時空を内在的にコンパクト化し、通常時空は非コンパクトのままである。

したがって双対時空回転角のみを離散化する。
\[
\theta \in \frac{2\pi}{N} \mathbb{Z},
\]
$N \gg 1$ は双対時空のコンパクト化スケールを決める非常に大きな整数（典型的には粒子の静止質量のプランク単位への比例）である。通常時空ラピディティは連続のままである。
\[
\phi \in \mathbb{R}.
\]

この非対称性は理論的一貫性のために必要である。
\begin{itemize}
  \item 双対時空は楕円的（$\hat{q}^2 = -1$）で内在的にコンパクトであり、その離散化は $R_{\text{dual}}$ によってパラメーター化される $S^3$ の $4\pi$ 周期的位相幾何を直接反映する。これはニュートン--ウィグナー局在化問題を解決し自然な量子離散性を与えるために不可欠な有限位相体積を提供する。
  \item 通常時空は双曲的（$\hat{p}^2 = +1$）で内在的に非コンパクトであるから、$\phi$ の連続性は完全ローレンツ群を保ち、巨視的極限における古典 GR との正確な等価性を保証する。
\end{itemize}

ねじれ不整合ローター $\Omega = R_{\text{usual}}^\dagger R_{\text{dual}}$ はしたがって連続的通常成分と離散的双対成分を結合する。ねじれスカラー $J = \frac{1}{16} B(\Omega_{\text{biv}}, \Omega_{\text{biv}})$ は有界のままであり、連続体誘起特異点を厳密に排除する。

ねじれ層を跨ぐ有効ローレンツ因子に典型的に現れる同調関数は、次の形をとる。
\[
\gamma_{\text{eff}} \propto D_N(\delta\theta) \equiv \frac{\sin(N \delta\theta / 2)}{N \sin(\delta\theta / 2)},
\]
ここで $\delta\theta$ はねじれ不整合角であり、$D_N$ は（正規化された）ディリクレ核である。分母は格子 $\delta\theta \in 2\pi \mathbb{Z}$ 上で正確に零になる。

そのような各点で分子 $\sin(N \cdot k\pi)$ も（$N$ が素数か合成かを問わず任意の整数に対して）零となるため、見かけの特異点は実際にはディリクレ核型の除去可能な $0/0$ である。標準的なロピタル評価により有限極限
\[
\lim_{\delta\theta \to 2\pi k} D_N(\delta\theta) = \pm 1 \qquad (k \in \mathbb{Z}),
\]
が得られ、同調関数はコンパクト格子全体で有界であり、古典連続体極限 $N \to \infty$ は病理なく回復される。

したがってねじれ動力学の有界性はディリクレ核の構造的性質であり、$N$ が素数である必要はない。任意の大きな正整数 $N$ で偽特異点を抑えつつ連続体極限を保てる。双対時空コンパクト化の数論的構造---特に $N$ と粒子質量スペクトルの関係---は枠組みの自然な拡張であり、離散双対時空モデルに関する姉妹論文で詳述される。

\section{双対ローター動力学とド・ブロイ位相の幾何的起源}
\label{sec:dualrotordynamics}

二重時空理論では、ねじれ不整合スカラー $J$ は相対ローター $\Omega = R_{\text{usual}}^\dagger R_{\text{dual}}$ から構成される。$J$ は古典場理論における重力作用密度として機能する一方、個々の質量粒子の内部自由度を支配するポテンシャルとしての自然な粒子内在的解釈も許す。

小さな不整合角に対して、キリング形式不変量は
\[
J \approx \frac{1}{2} \sum_{a=1}^3 \delta\phi_a^2,
\]
と展開する。ここで $\delta\phi_a = \phi_a - \theta_a$ は各ブースト--回転平面におけるねじれ不整合角である。双対セクターの剛性定数として静止質量 $m$ を組み込むと
\[
J = \frac{1}{2} m \sum_{a=1}^3 \delta\phi_a^2.
\]

単一粒子的内在ローター角に対する有効作用を提案する。
\[
S_{\text{particle}} = \int \left[ \frac{1}{2} \sum_{a=1}^3 \dot{\phi}_a^2 - \frac{1}{2} \sum_{a=1}^3 \dot{\theta}_a^2 - \frac{m}{2} \sum_{a=1}^3 (\phi_a - \theta_a)^2 \right] dt.
\]
運動項の相反符号は通常セクターのブースト様性質と双対セクターの時間反転回転様性質を反映する。

オイラー--ラグランジュ方程式は
\begin{align}
\ddot{\phi}_a &= m (\phi_a - \theta_a), \\
-\ddot{\theta}_a &= m (\phi_a - \theta_a).
\end{align}
ねじれ不整合角 $\delta\phi_a = \phi_a - \theta_a$ を定義すると、単調振動子方程式
\begin{equation}
\ddot{\delta\phi}_a + m \delta\phi_a = 0.
\label{eq:dual-harmonic}
\end{equation}
を得る。振動周波数は $\sqrt{m}$（$\hbar = c = 1$ の自然単位）である。

一様運動の自由粒子では、通常ラピディティ成分は一定率 $\dot{\phi}_a = \gamma v_a/c$ で進む。式~\eqref{eq:dual-harmonic}の一般解は
\[
\delta\phi_a(t) = A_a \sin(\sqrt{m}\, t + \psi_a).
\]
低エネルギー非相対論的領域では、速いコンプトン周波数振動は平均して零になり、固有時に比例する位相の遅い線形蓄積が残る。$\hbar$ と $c$ を復元すると、蓄積位相は
\[
\int \delta\phi_a \, dt \propto \frac{E t - \mathbf{p} \cdot \mathbf{x}}{\hbar},
\]
となり、これはまさに物質波のド・ブロイ位相である。

したがって量子力学波動関数位相 $S(x,t)/\hbar$ は、通常ローターに対する双対ローターの遅れの緩やかに変化する成分として幾何学的に創発する。これは二重時空枠組みの内で、付加公準なしに量子物質波の完全に背景独立な代数的起源を与える。

この粒子内在的作用の多粒子への集合平均は古典テレパラレル作用 $\int J \, d^4x$ を回復し、巨視的スケールで一般相対性理論との完全な一貫性を保証しつつ、双対時空自由度の微視的量子動力学を明らかにする。

この調和振動子動学は、重力の古典ねじれ記述と次節で展開する完全量子理論との橋である。

\section{量子重力と完全統一に向けて}
\label{sec:quantumgravity}

ここまでの古典定式化は、粒子内在的通常・双対時空間のねじれ不整合として重力が創発すること、それが $\mathrm{Cl}(3,1)$ と同型の16次元双四元数代数に符号化されることを示した。作用 $S$ はアインシュタイン--ヒルベルト作用と正確に等価でありながら、連続体仮説とリーマン曲率を捨てる。

深い拡張は、双対ローター $R_\text{dual}$ が数学的に $\mathrm{SU}(2)$ スピン1/2回転演算子と同一であることの認識から生じる。したがって双対セクター位相 $\theta$ は量子状態を直接パラメーター化する。双対時空の内在的コンパクト性（$\theta \bmod 4\pi$）により、これらの位相はコンパクト多様体上にあり、各粒子の量子状態 $|\psi\rangle$ が $R_\text{dual}$ のユニタリ位相部分に対応するヒルベルト空間構造を自然に与える。

位置と運動量演算子は双対ローター角の射影として創発し、重力自身を含む力は相対ローター位相の変調として現れる。これは量子力学を双対時空回転と、重力を同一代数的枠組みにおけるねじれ不整合と同定し、有界離散ねじれ層を通じて特異点を排除する。

明示的ヒルベルト空間構成、$\mathrm{Cl}(3,1)$ の最小左イデアルへのスピノル埋め込み、すべての相互作用の統一、およびプランク期時間矢固定を通じたバリオン非対称の解決を含む完全量子重力枠組みは、姉妹論文「双四元数二重時空におけるねじれ不整合としての量子重力」\cite{dst-quantum-gravity}で詳述される。

この拡張は理論を完成させる。一般相対性理論は古典的に保存される一方、量子重力は余次元やアドホック量子化なしに、単に各質量粒子に内在的な二重時空構造から代数的に創発する。

\section{結論}
\label{sec:conclusions}
二重時空理論は、双四元数代数 $\mathbb{H} \otimes_{\mathbb{R}} \mathbb{H} \cong \mathrm{Cl}(3,1)$ に基づき、重力、慣性、および基本相互作用を理解するための変革的枠組みを提供する。各粒子を絡み合った一対のミンコフスキー時空（通常と双対）の中に埋め込むことで、理論は重力現象をこれらのセクター間のねじれ的ずれの現れとして再解釈する。このパラダイム転換は重力物理学の長年のパラドックスを解決し、層状ねじれ動学の現れとして強い核力を統一し、内在的カイラリティと時間矢の固定を通じたバリオン非対称の自然な機構を提供する。双対ローター同期を通じた重力制御から、原子核を原始ねじれ星として再解釈するまで、理論の予測は実験的検証と技術革新の新たな道を開く。重力物理学の新時代の入口に立つ今、双対時空枠組みは宇宙と我々のその中における位置についての理解を再考するよう招く。

\bibliographystyle{unsrt}
\begin{thebibliography}{1}

\bibitem{einstein1928}
A.~Einstein,
``Riemann-Geometrie mit Aufrechterhaltung des Begriffes des Fernparallelismus,''
\textit{Sitzungsber. Preuss. Akad. Wiss. Phys.-Math. Kl.},
217--221 (1928).

\bibitem{einstein1939}
A.~Einstein,
``On a stationary system with spherical symmetry consisting of many gravitating masses,''
\textit{Ann. Math.},
\textbf{40}, no.~4, 922--936 (1939).
doi:10.2307/1968902

\bibitem{girard1999}
P.~R. Girard,
``Einstein’s equations and Clifford algebra,''
\textit{Advances in Applied Clifford Algebras},
\textbf{9}, no.~2, 225--230 (1999).
doi:10.1007/BF03042377

\bibitem{maluf2013}
J.~W. Maluf,
``The teleparallel equivalent of general relativity,''
\textit{Ann. Phys. (Berlin)},
\textbf{525}, 339--359 (2013).
doi:10.1002/andp.201200272

\bibitem{dst-quantum-gravity} 匿名 2025, 「双四元数二重時空におけるねじれ不整合としての量子重力」, https://github.com/hypernumbernet/dual-spacetime-doc/blob/main/dst-quantum-gravity.pdf.

\end{thebibliography}

\clearpage
\appendix
\section{乗法表}
\label{app:mult-table}

\begin{table}[ht]
  \begin{center}
    \caption{双四元数乗法表}
    \small
    \begin{tabular}{r|rrrrrrrrrrrrrrr}
      &$j$&$kI$&$kJ$&$kK$&$iI$&$iJ$&$iK$&$I$&$J$&$K$&$k$&$jI$&$jJ$&$jK$&$i$ \\ \hline
      j&$-1$&$+iI$&$+iJ$&$+iK$&$-kI$&$-kJ$&$-kK$&$+jI$&$+jJ$&$+jK$&$+i$&$-I$&$-J$&$-K$&$-k$ \\ \hline
      kI&$-iI$&$+1$&$-K$&$+J$&$-j$&$+jK$&$-jJ$&$-k$&$+kK$&$-kJ$&$-I$&$+i$&$-iK$&$+iJ$&$+jI$ \\ \hline
      kJ&$-iJ$&$+K$&$+1$&$-I$&$-jK$&$-j$&$+jI$&$-kK$&$-k$&$+kI$&$-J$&$+iK$&$+i$&$-iI$&$+jJ$ \\ \hline
      kK&$-iK$&$-J$&$+I$&$+1$&$+jJ$&$-jI$&$-j$&$+kJ$&$-kI$&$-k$&$-K$&$-iJ$&$+iI$&$+i$&$+jK$ \\ \hline
      iI&$+kI$&$+j$&$-jK$&$+jJ$&$+1$&$-K$&$+J$&$-i$&$+iK$&$-iJ$&$-jI$&$-k$&$+kK$&$-kJ$&$-I$ \\ \hline
      iJ&$+kJ$&$+jK$&$+j$&$-jI$&$+K$&$+1$&$-I$&$-iK$&$-i$&$+iI$&$-jJ$&$-kK$&$-k$&$+kI$&$-J$ \\ \hline
      iK&$+kK$&$-jJ$&$+jI$&$+j$&$-J$&$+I$&$+1$&$+iJ$&$-iI$&$-i$&$-jK$&$+kJ$&$-kI$&$-k$&$-K$ \\ \hline
      I&$+jI$&$-k$&$+kK$&$-kJ$&$-i$&$+iK$&$-iJ$&$-1$&$+K$&$-J$&$+kI$&$-j$&$+jK$&$-jJ$&$+iI$ \\ \hline
      J&$+jJ$&$-kK$&$-k$&$+kI$&$-iK$&$-i$&$+iI$&$-K$&$-1$&$+I$&$+kJ$&$-jK$&$-j$&$+jI$&$+iJ$ \\ \hline
      K&$+jK$&$+kJ$&$-kI$&$-k$&$+iJ$&$-iI$&$-i$&$+J$&$-I$&$-1$&$+kK$&$+jJ$&$-jI$&$-j$&$+iK$ \\ \hline
      k&$-i$&$-I$&$-J$&$-K$&$+jI$&$+jJ$&$+jK$&$+kI$&$+kJ$&$+kK$&$-1$&$-iI$&$-iJ$&$-iK$&$+j$ \\ \hline
      jI&$-I$&$-i$&$+iK$&$-iJ$&$+k$&$-kK$&$+kJ$&$-j$&$+jK$&$-jJ$&$+iI$&$+1$&$-K$&$+J$&$-kI$ \\ \hline
      jJ&$-J$&$-iK$&$-i$&$+iI$&$+kK$&$+k$&$-kI$&$-jK$&$-j$&$+jI$&$+iJ$&$+K$&$+1$&$-I$&$-kJ$ \\ \hline
      jK&$-K$&$+iJ$&$-iI$&$-i$&$-kJ$&$+kI$&$+k$&$+jJ$&$-jI$&$-j$&$+iK$&$-J$&$+I$&$+1$&$-kK$ \\ \hline
      i&$+k$&$-jI$&$-jJ$&$-jK$&$-I$&$-J$&$-K$&$+iI$&$+iJ$&$+iK$&$-j$&$+kI$&$+kJ$&$+kK$&$-1$ \\ \hline
    \end{tabular}
  \end{center}
\end{table}
\begin{table}[ht]
  \begin{center}
    \caption{$\mathrm{Cl}(3,1)$ 変形乗法表（$e_0, e_1, e_2, e_3 \mapsto 0,1,2,3$）}
    \scriptsize
    \begin{tabular}{r|rrrrrrrrrrrrrrr}
      &$0$&$1$&$2$&$3$&$01$&$02$&$03$&$32$&$13$&$21$&$123$&$032$&$013$&$021$&$0123$\\\hline
      $0$&$-$&$+01$&$+02$&$+03$&$-1$&$-2$&$-3$&$+032$&$+013$&$+021$&$+0123$&$-32$&$-13$&$-21$&$-123$\\
      $1$&$-01$&$+$&$-21$&$+13$&$-0$&$+021$&$-013$&$-123$&$+3$&$-2$&$-32$&$+0123$&$-03$&$+02$&$+032$\\
      $2$&$-02$&$+21$&$+$&$-32$&$-021$&$-0$&$+032$&$-3$&$-123$&$+1$&$-13$&$+03$&$+0123$&$-01$&$+013$\\
      $3$&$-03$&$-13$&$+32$&$+$&$+013$&$-032$&$-0$&$+2$&$-1$&$-123$&$-21$&$-02$&$+01$&$+0123$&$+021$\\
      $01$&$+1$&$+0$&$-021$&$+013$&$+$&$-21$&$+13$&$-0123$&$+03$&$-02$&$-032$&$-123$&$+3$&$-2$&$-32$\\
      $02$&$+2$&$+021$&$+0$&$-032$&$+21$&$+$&$-32$&$-03$&$-0123$&$+01$&$-013$&$-3$&$-123$&$+1$&$-13$\\
      $03$&$+3$&$-013$&$+032$&$+0$&$-13$&$+32$&$+$&$+02$&$-01$&$-0123$&$-021$&$+2$&$-1$&$-123$&$-21$\\
      $32$&$+032$&$-123$&$+3$&$-2$&$-0123$&$+03$&$-02$&$-$&$+21$&$-13$&$+1$&$-0$&$+021$&$-013$&$+01$\\
      $13$&$+013$&$-3$&$-123$&$+1$&$-03$&$-0123$&$+01$&$-21$&$-$&$+32$&$+2$&$-021$&$-0$&$+032$&$+02$\\
      $21$&$+021$&$+2$&$-1$&$-123$&$+02$&$-01$&$-0123$&$+13$&$-32$&$-$&$+3$&$+013$&$-032$&$-0$&$+03$\\
      $123$&$-0123$&$-32$&$-13$&$-21$&$+032$&$+013$&$+021$&$+1$&$+2$&$+3$&$-$&$-01$&$-02$&$-03$&$+0$\\
      $032$&$-32$&$-0123$&$+03$&$-02$&$+123$&$-3$&$+2$&$-0$&$+021$&$-013$&$+01$&$+$&$-21$&$+13$&$-1$\\
      $013$&$-13$&$-03$&$-0123$&$+01$&$+3$&$+123$&$-1$&$-021$&$-0$&$+032$&$+02$&$+21$&$+$&$-32$&$-2$\\
      $021$&$-21$&$+02$&$-01$&$-0123$&$-2$&$+1$&$+123$&$+013$&$-032$&$-0$&$+03$&$-13$&$+32$&$+$&$-3$\\
      $0123$&$+123$&$-032$&$-013$&$-021$&$-32$&$-13$&$-21$&$+01$&$+02$&$+03$&$-0$&$+1$&$+2$&$+3$&$-$\\
    \end{tabular}
  \end{center}
\end{table}
これらは同一の結果であり同型である。
\clearpage
\section{ねじれ不整合スカラー J と TEGR のテレパラレルねじれスカラー T の等価性}
\label{app:TEGR-equivalence}

ねじれ不整合バイベクトルは
\[
\Omega_{\rm biv} = \log(R_{\rm usual}^\dagger R_{\rm dual}) = \sum_{a=1}^3 \left( \frac{\alpha_a}{2} i\Gamma_a + \frac{\beta_a}{2} \Gamma_a \right),
\]
であり、\( (i\Gamma_a)^2 = +1 \), \( \Gamma_a^2 = -1 \), \([i\Gamma_a, \Gamma_b] = 0\) である。

\(\mathfrak{so}(3,1) \oplus \mathfrak{so}(3,1)\) 上のキリング形式は \( B(i\Gamma_a, i\Gamma_a) = +8 \), \( B(\Gamma_a, \Gamma_a) = -8 \) を満たし、
\[
B(\Omega_{\rm biv}, \Omega_{\rm biv}) = 8 \sum_{a=1}^3 (\alpha_a^2 - \beta_a^2).
\]
したがってねじれ不整合スカラーは
\[
J = \frac{1}{16} B(\Omega_{\rm biv}, \Omega_{\rm biv}) = \frac12 \sum_{a=1}^3 (\alpha_a^2 - \beta_a^2).
\]

この不変量を TEGR と比較するには、まずローター場とテトラッドの幾何的対応を指定しなければならない。具体的には、全ローター場 $R_{\rm total}(x) = R_{\rm usual}(x)\, R_{\rm dual}(x)$ が固定参照テトラッド $\{\mathring{e}^a\}$ にスピノルサンドイッチ積で作用する。
\[
e^a{}_\mu(x) \;=\; \langle R_{\rm total}(x)\, \mathring{e}^a\, R_{\rm total}^\dagger(x) \rangle_{1,\mu},
\]
ここで $\langle \cdot \rangle_{1,\mu}$ は座標基底に沿った1次（ベクトル）成分を抽出する。この同定により $R_{\rm total}$ は計量 $g_{\mu\nu} = \eta_{ab}\, e^a{}_\mu e^b{}_\nu$ をもつテトラッド場 $e^a{}_\mu$ へと昇格し、双対不整合バイベクトル $\Omega_{\rm biv}$ は $e^a{}_\mu$ に付随する Weitzenböck 接続の contorsion に比例する。
\[
\Omega_{\rm biv}\big|_{\mu\nu}^{ab} \;\propto\; K^{ab}{}_\mu\, dx^\mu \wedge dx^\nu,
\qquad
T^a{}_{\mu\nu} \;=\; \partial_\mu e^a{}_\nu - \partial_\nu e^a{}_\mu.
\]
この同定の下で二次不変量 $J$ は（全発散まで）テレパラレルねじれスカラー
\[
T = \frac14 T_{\rho\mu\nu}T^{\rho\mu\nu} + \frac12 T_{\rho\mu\nu}T^{\nu\mu\rho} - T_\rho T^\rho,
\]
に対応し、オンシェルでは
\[
J = \frac12 T + \nabla_\mu B^\mu,
\]
となる。ここで $\nabla_\mu B^\mu$ は上の射影によって形が定まる全発散である。

TEGR の標準的結果 \cite{maluf2013} により、同一のテトラッドに対して
\[
\int T\, e\, d^4x = -\int R\sqrt{-g}\, d^4x + \int \partial_\mu(\sqrt{-g}\, V^\mu)\, d^4x,
\]
すなわち $T$ と $-R$ は全発散のみ異なる。$J$ の関係を代入すると
\[
S = \frac{c^4}{16\pi G} \int J \, d^4x = -\frac{c^4}{16\pi G} \int R \sqrt{-g}\, d^4x + (\text{境界項}).
\]

$S$ の変分は、テトラッド $e^a{}_\mu$（等価的にはローター場 $R_{\rm total}$）を独立変数とし、積分領域の境界 $\partial\mathcal{M}$ 上ですべての場の変分 $\delta e^a{}_\mu$ が零になる標準的境界条件の下で行われる。この条件下では上の全発散項は変分方程式に寄与せず、残りの変分はアインシュタイン方程式
\[
G_{\mu\nu} = \frac{8\pi G}{c^4} T_{\mu\nu}.
\]
を再現する。

したがって、上で述べたローター--テトラッド同定と標準的変分境界条件の下で、二重時空作用は古典極限においてアインシュタイン--ヒルベルト作用と動学的に等価である。

\end{document}
